\documentclass[11pt,a4paper]{article}
% =====================================================================
% Préambule commun au rapport (compilation pdfLaTeX + biber, Overleaf)
% =====================================================================
\usepackage[T1]{fontenc}
\usepackage[utf8]{inputenc}
\usepackage{lmodern}
\usepackage[french]{babel}
\usepackage{microtype}
\usepackage[a4paper,margin=2cm]{geometry}

\usepackage{amsmath,amssymb,mathtools}
\usepackage[locale=FR]{siunitx}   % virgule décimale, espaces fines
\sisetup{exponent-product=\times, per-mode=symbol}
% locale=FR ne traduit pas les mots des intervalles et des listes (« to », « and »)
\sisetup{range-phrase={~à~}, list-pair-separator={ et }, list-final-separator={ et },
  list-separator={, }}

\usepackage{graphicx}
\usepackage{subcaption}
\usepackage{booktabs}
\usepackage{tabularx}
\usepackage{longtable}
\usepackage{xcolor}
\usepackage{enumitem}
\setlist{nosep}

\usepackage[backend=biber,style=numeric,sorting=nyt]{biblatex}
\addbibresource{biblio.bib}
\renewcommand*{\bibfont}{\small}

\usepackage[hidelinks,hypertexnames=false]{hyperref}
\usepackage[french,noabbrev]{cleveref}  % toujours après hyperref

% =====================================================================
% Macros : notations de l'énoncé et cadres « à venir »
% =====================================================================

% ---------- Notations (identiques à l'énoncé, au code et aux slides) ----------
\newcommand{\Fe}{F_e}
\newcommand{\Nfft}{N_{\mathrm{fft}}}
\newcommand{\Tw}{T_w}
\newcommand{\Tq}{T_q}
\newcommand{\SdB}{S_{\mathrm{dB}}}
\newcommand{\sthr}{s_{\mathrm{thr}}}
\newcommand{\gf}{\gamma_f}
\newcommand{\Nf}{N_f}
\newcommand{\Nt}{N_t}
\newcommand{\Nu}{N_u}
\newcommand{\nq}{n_q}
\newcommand{\Meff}{M_{\mathrm{eff}}}
\newcommand{\kw}{\kappa_w}
\newcommand{\mv}{m_v}
\newcommand{\kv}{k_v}
\newcommand{\dm}{\Delta m}
\newcommand{\dk}{\Delta k}
\newcommand{\RSB}{\ensuremath{\mathrm{RSB}}}
\newcommand{\TBD}{\ensuremath{\mathrm{TBD}}}
\newcommand{\TFD}{\ensuremath{\mathrm{TFD}}}
\newcommand{\hhat}{\hat{h}}
\newcommand{\E}{\mathbb{E}}
\newcommand{\Prob}{\mathbb{P}}

% ---------- Texte à compléter (rouge, bien visible) ----------
% Avant le rendu : chercher « À remplir » dans le PDF, il ne doit plus y en avoir.
\DeclareRobustCommand{\aremplir}[1]{\textcolor{red}{\textbf{[À remplir : #1]}}}

% ---------- Figures de A : cadre « figure à venir » tant que le PDF n'existe pas ----------
% Usage : \inclurefig{tbd_rsb}{ce que la figure montrera}
% Cherche <nom>.pdf puis <nom>.png, dans figures/ (Overleaf) puis ../figures/ (dépôt git).
% La figure est insérée à sa taille naturelle : MATLAB l'exporte déjà à sa largeur
% d'impression (utils/sauver_figure.m), donc son texte reste à 8-9 pt. Elle n'est
% réduite que si elle dépasse la largeur disponible. L'argument optionnel est ignoré
% (gardé pour compatibilité).
\newsavebox{\boitefig}
\newcommand{\figurenaturelle}[1]{%
  \sbox{\boitefig}{\includegraphics{#1}}%
  \ifdim\wd\boitefig>\linewidth\resizebox{\linewidth}{!}{\usebox{\boitefig}}\else\usebox{\boitefig}\fi}
\newcommand{\inclurefig}[3][]{%
  \IfFileExists{figures/#2.pdf}{\figurenaturelle{figures/#2.pdf}}{%
  \IfFileExists{figures/#2.png}{\figurenaturelle{figures/#2.png}}{%
  \IfFileExists{../figures/#2.pdf}{\figurenaturelle{../figures/#2.pdf}}{%
  \IfFileExists{../figures/#2.png}{\figurenaturelle{../figures/#2.png}}{%
  \fbox{\parbox[c][3.2cm][c]{0.92\linewidth}{\centering\small\color{gray}
    Figure à venir : #3\\[2pt] \texttt{\detokenize{#2.pdf}}}}}}}}}

% ---------- Tableaux de A (fichiers .tex booktabs, sans environnement table) ----------
% Usage : \inclureTableau{T02}{ce que le tableau contiendra}
\newcommand{\inclureTableau}[2]{%
  \IfFileExists{tableaux/#1.tex}{\input{tableaux/#1}}{%
  \IfFileExists{../resultats/tableaux/#1.tex}{\input{../resultats/tableaux/#1}}{%
  \fbox{\parbox[c][1.6cm][c]{0.92\linewidth}{\centering\small\color{gray}
    Tableau à venir : #2\\[2pt] \texttt{\detokenize{#1.tex}}}}}}}

% Case de tableau à remplir (version courte, pour ne pas élargir les tableaux)
\DeclareRobustCommand{\aq}{\textcolor{red}{\textbf{?}}}



\title{Reconnaissance musicale par empreintes audio\\[4pt]
  \large Comment un téléphone reconnaît-il un morceau en 5 secondes ?}
\author{Alessandro \aremplir{nom} \and Laure Dauthieu \and \aremplir{3\textsuperscript{e} membre du trinôme ?}\\[2pt]
  \small ESIEE Paris -- OBL-4101 Traitement du signal -- 2026--2027}
\date{27 novembre 2026}

\begin{document}
\maketitle
% Les équations (1) à (8) de l'énoncé sont repérées par \tag ; les nôtres partent de (9).
\setcounter{equation}{8}
% ---------------------------------------------------------------------
% Résumé (5 lignes, chiffres clés) + où trouver les réponses Q1 à Q6
% ---------------------------------------------------------------------
\begin{abstract}
\noindent Nous identifions un morceau à partir d'un extrait court, bruité et pris à un
instant inconnu, en suivant la méthode de Wang : constellation de maxima locaux du
spectrogramme, clés formées de paires de pics, index inversé et vote sur le décalage
temporel. Sur une base de \aremplir{100} morceaux et \aremplir{x} requêtes positives et
\aremplir{y} négatives, le système atteint un \TBD{} de \aremplir{x~\%} pour un \TFD{} de
\aremplir{y~\%} ($\Tq = \aremplir{5}$\,s, $\sthr = \aremplir{s}$, bruit blanc à
\qty{10}{\decibel}). La théorie explique pourquoi l'empreinte est invariante au gain et au
filtrage, prévoit un score en $F\rho\Tq p^2$ et fixe un plancher au seuil de décision ;
\aremplir{phrase sur la confrontation à la mesure}.
\end{abstract}

\begin{table}[h]
  \centering\small
  \caption{Où trouver les réponses aux questions de l'énoncé.}
  \label{tab:ou-reponses}
  \begin{tabularx}{\linewidth}{@{}lXl@{}}
    \toprule
    Question & Sujet & Où \\
    \midrule
    Q1 & Choix de $\Fe = \qty{8}{\kilo\hertz}$ : Shannon, ce que l'on perd
       & \cref{sec:shannon} \\
    Q2 & Tableau $\Tw$, $\Delta f$, trames/s ; choix de $N$ ; propriété (2) démontrée et vérifiée
       & \cref{sec:tfct,sec:proprietes}, \cref{tab:resolutions,fig:filtrage} \\
    Q3 & Voisinage, densité $\rho$, survie $p$ ; pourquoi $\rho p^2$
       & \cref{sec:pics}, \cref{fig:p-rho} \\
    Q4 & $p$ mesuré, tableau pic/paire ; choix de $\gf$
       & \cref{sec:empreintes}, \cref{tab:pic-paire,tab:gammaf} \\
    Q5 & Histogrammes des scores ; seuil automatique
       & \cref{sec:decision}, \cref{fig:histo-scores,fig:tbd-tfd} \\
    Q6 & MP3 contre bruit blanc ; genres difficiles
       & \cref{sec:q6,sec:genres}, \cref{fig:genres} \\
    \bottomrule
  \end{tabularx}
\end{table}


\tableofcontents
\clearpage

% ---------------------------------------------------------------------
\section{Introduction}
\label{sec:intro}
% ---------------------------------------------------------------------

Reconnaître un morceau à partir de quelques secondes captées par un téléphone, dans un
bar bruyant, à un instant quelconque du morceau, demande une représentation du signal
qui concilie deux exigences opposées. Elle doit être \emph{robuste} : insensible au
volume, au haut-parleur, à la salle, au microphone, au codage MP3, au bruit ambiant et à
la position de l'extrait. Elle doit aussi être \emph{spécifique} : deux morceaux différents
ne doivent presque jamais produire la même empreinte.

Deux familles de solutions se sont imposées au début des années 2000. Haitsma et
Kalker~\cite{haitsma2002} découpent le spectre en 33~bandes logarithmiques entre 300 et
\qty{2000}{\hertz} et codent, trame après trame, le signe des différences d'énergie entre
bandes voisines et trames successives : on obtient une suite de mots de 32~bits, et deux
extraits sont comparés par leur taux d'erreur binaire. Wang~\cite{wang2003}, à l'origine
de Shazam, ne garde au contraire que les maxima locaux du spectrogramme, une
\emph{carte de constellation}, et forme des clés à partir de \emph{paires} de pics : la clé
code les deux fréquences et l'écart de temps, l'instant absolu étant rangé à part. Une
requête vote alors pour un décalage temporel, et le bon morceau se signale par une
accumulation de votes sur un seul décalage. Ces deux approches partagent la même idée :
ne retenir que des comparaisons \emph{relatives} entre valeurs voisines du plan
temps-fréquence, qui survivent au gain et au filtrage.

Nous avons suivi la méthode de Wang, implémentée sous MATLAB et évaluée sur
\aremplir{100} morceaux de la base libre FMA~\cite{fma2017}. Au-delà d'une chaîne qui
fonctionne, nous avons cherché à \emph{prédire} les performances : le score attendu du bon
morceau, le plancher du seuil de décision et l'effet de chaque dégradation sur les pics,
puis à confronter ces prédictions à la mesure. \aremplir{Résultat principal en une
phrase, par exemple : « à $\Tq = \qty{5}{\second}$, le système reconnaît x~\% des extraits
jusqu'à y~dB de RSB, et la théorie prévoyait la rupture à z~dB. »}

La \cref{sec:modelisation} décrit le modèle de dégradation, la \cref{sec:theorie} démontre
les propriétés d'invariance sur lesquelles repose la méthode, la \cref{sec:conception}
justifie chaque étape et chaque paramètre, la \cref{sec:resultats} présente le protocole et
les mesures, et la \cref{sec:discussion} analyse les échecs. Les notations sont celles de
l'énoncé (\cref{sec:notations}).

% ---------------------------------------------------------------------
\section{Modélisation}
\label{sec:modelisation}
% ---------------------------------------------------------------------

Entre le fichier d'origine $x(t)$ et l'extrait capté $y(t)$ s'intercalent un haut-parleur,
une salle, un microphone, un codec et du bruit ambiant. L'énoncé les résume par
\begin{equation}
  y(t) = \alpha\,(h * x)(t - t_0) + b(t),
  \tag{1}\label{eq:modele}
\end{equation}
auquel s'ajoute le codage perceptuel, qui n'est pas linéaire et ne s'écrit pas sous cette
forme. Chaque terme a une cause physique et appelle une invariance précise de l'empreinte
(\cref{tab:modele}).

\begin{itemize}
  \item $\alpha > 0$ est un \textbf{gain} inconnu : volume de lecture, distance à la source,
        sensibilité du micro. Il est constant à l'échelle d'un extrait.
  \item $h$ est la \textbf{réponse impulsionnelle} de la chaîne haut-parleur $\to$ salle
        $\to$ micro. Sa partie courte (haut-parleur, micro, premières réflexions) colore le
        spectre ; sa partie longue (réverbération tardive) étale le signal dans le temps.
  \item $t_0$ est le \textbf{décalage} inconnu : on ignore où l'extrait a été capté.
  \item $b(t)$ est un \textbf{bruit additif} indépendant de $x$ : souffle blanc, brouhaha
        coloré et non stationnaire.
  \item Le \textbf{codage MP3} supprime les composantes masquées et ajoute un bruit de
        quantification mis en forme sous le seuil de masquage.
\end{itemize}

\begin{table}[h]
  \centering\small
  \caption{Terme du modèle~\eqref{eq:modele}, dégradation qui le simule et invariance visée.}
  \label{tab:modele}
  \begin{tabularx}{\linewidth}{@{}lXX@{}}
    \toprule
    Terme & Dégradation simulée & Invariance visée (où) \\
    \midrule
    $t_0$      & début de requête tiré au hasard, non aligné sur la grille
               & translation du spectrogramme, vote sur $\delta$ (\cref{sec:proprietes,sec:decision}) \\
    $\alpha$   & implicite dans toutes les conditions
               & positions des maxima, seuil local (\cref{sec:maxima}) \\
    $h$ court  & passe-bande \qtyrange{300}{3400}{\hertz}, haut-parleur de téléphone
               & biais $g(k)$ indépendant du temps, condition (5) (\cref{sec:maxima}) \\
    $h$ long   & réverbération $T_{60} = \qty{0.3}{\second}$ et \qty{1}{\second}
               & hors modèle : limite de la propriété (2) (\cref{sec:limites}) \\
    $b$        & bruit blanc (\qtyrange{20}{0}{\decibel}), brouhaha (\qtyrange{10}{0}{\decibel})
               & survie des pics, RSB local par bande (\cref{sec:q6}) \\
    MP3        & \num{128}, \num{64} et \qty{32}{\kilo\bit\per\second}
               & les pics sont les masqueurs (\cref{sec:q6}) \\
    tout       & enregistrements au micro en salle, à deux distances
               & toutes à la fois (\cref{sec:reel}) \\
    \bottomrule
  \end{tabularx}
\end{table}

Les deux exigences sont antagonistes : une empreinte riche (beaucoup de pics, clés fines)
est spécifique mais fragile, une empreinte pauvre est robuste mais confond les morceaux.
Toute la conception (\cref{sec:conception}) consiste à quantifier ce compromis : la
robustesse par la probabilité $p$ qu'un pic survive, la spécificité par le nombre de clés
possibles $M$ et par le nombre d'appariements fortuits qu'il entraîne.

% ---------------------------------------------------------------------
\section{Théorie}
\label{sec-theorie}
% ---------------------------------------------------------------------

Nous gardons la numérotation de l'énoncé : les équations (1) à (8) et les propriétés (1)
(translation) et (2) (filtrage) y renvoient directement. Nos propres équations sont
numérotées à partir de (9).

\subsection{Prétraitement et théorème de Shannon (Q1)}
\label{sec-shannon}

Le signal est ramené en mono (moyenne des canaux) puis rééchantillonné à
$\Fe = \qty{8}{\kilo\hertz}$.

\paragraph{Ce qu'impose Shannon.} Un signal échantillonné à $\Fe$ n'est représenté sans
ambiguïté que s'il ne contient aucune composante au-delà de $\Fe/2$ : une sinusoïde de
fréquence $f_0 \in\, ]\Fe/2, \Fe[$ donne exactement les mêmes échantillons qu'une
sinusoïde de fréquence $\Fe - f_0$. Choisir $\Fe = \qty{8}{\kilo\hertz}$, c'est donc décider
que seule la bande $[0, \qty{4}{\kilo\hertz}]$ servira, et s'obliger à éliminer tout ce qui
est au-dessus \emph{avant} de décimer.

\paragraph{Pourquoi la bande \qtyrange{0}{4}{\kilo\hertz} suffit.} L'identité d'un morceau
tient à sa mélodie, à son harmonie et à son rythme, c'est-à-dire aux fondamentales des
notes et à leurs premiers harmoniques. La fondamentale la plus aiguë du piano
(do\textsubscript{8}) vaut \qty{4186}{\hertz} ; celles des voix chantées et de la quasi-totalité
des instruments mélodiques restent sous \qty{2}{\kilo\hertz}. L'essentiel de l'énergie
musicale est concentré dans cette zone, et c'est là que se trouvent les maxima les plus
marqués du spectrogramme, donc les pics les plus stables. La bande téléphonique
(\qtyrange{300}{3400}{\hertz}) a été dimensionnée sur le même constat : elle suffit à
reconnaître une voix ou une mélodie.

\paragraph{Ce que l'on perd.} On perd la \og brillance \fg{} : harmoniques élevés,
cymbales, attaques, sifflantes. Trois raisons rendent cette perte acceptable :
\begin{enumerate}
  \item ce contenu est peu énergétique, donc il porte peu de pics, et des pics fragiles ;
  \item c'est celui que le canal détruit en premier : un petit haut-parleur coupe vers
        \qty{3.4}{\kilo\hertz}, l'absorption de l'air et des parois croît avec la fréquence,
        et le codeur MP3 supprime les hautes fréquences dès que le débit baisse. Une
        empreinte qui s'appuierait sur ces fréquences serait présente dans la référence et
        absente de la requête ;
  \item le coût diminue : le nombre d'échantillons est divisé par
        $\num{44100}/\num{8000} \approx \num{5.5}$, comme le nombre de canaux $K$ à
        résolution temporelle égale. Avec le facteur $\log_2 N$ de la FFT ($N$ passerait de
        1024 à environ 5600 pour garder $\Tw = \qty{128}{\milli\second}$), le coût de la
        TFCT baisse d'un facteur de l'ordre de 7.
\end{enumerate}

\paragraph{Pourquoi \texttt{resample} et pas une décimation naïve.} Garder un échantillon
sur $k$ sans filtrer replierait toute la bande \qtyrange{4}{22}{\kilo\hertz} dans
\qtyrange{0}{4}{\kilo\hertz}, en créant de faux pics. Le plus grave est que ce contenu
n'est pas le même dans la référence (fichier d'origine, aigus intacts) et dans la requête
(aigus coupés par le canal ou le MP3) : les faux pics différeraient des deux côtés et
détruiraient l'appariement sans cause visible. La fonction \texttt{resample} applique un
filtre passe-bas anti-repliement avant de décimer, ce qui supprime le problème.

\subsection{Transformée de Fourier à court terme et compromis temps-fréquence (Q2)}
\label{sec-tfct}

Le contenu musical évolue en permanence : une TFD calculée sur tout le morceau
mélangerait toutes les notes et ne dirait rien de leur succession. On calcule donc la
transformée de Fourier à court terme et le spectrogramme en décibels
\begin{equation}
  X[m,k] = \sum_{n=0}^{N-1} x[n+mH]\, w[n]\, e^{-j2\pi kn/\Nfft},
  \qquad
  \SdB[m,k] = 10\log_{10}\bigl(|X[m,k]|^2 + \varepsilon\bigr),
  \tag{2}
\end{equation}
avec une fenêtre de Hann périodique, $H = N/2$ et $\Nfft = N$. La trame $m$ et le canal $k$
correspondent à $t_m = mH/\Fe$ et $f_k = k\Fe/\Nfft$.

\paragraph{Résolutions.} La résolution temporelle est la durée de la fenêtre,
$\Tw = N/\Fe$. La résolution fréquentielle est la largeur à \qty{-3}{\decibel} du lobe
principal de la fenêtre, $\Delta f = \kw\,\Fe/N$, où $\kw$ ne dépend que de la forme de la
fenêtre. L'énoncé donne $\kw = \num{1.30}$, qui est la valeur de la fenêtre de
\emph{Hamming} ; pour la fenêtre de \emph{Hann} que nous utilisons, $\kw = \num{1.44}$
\cite{harris1978}. Le produit
\begin{equation}
  \Delta f \cdot \Tw = \kw
\end{equation}
ne dépend pas de $N$ : toute précision gagnée en fréquence est perdue en temps, et
inversement. Monter d'un demi-ton multiplie la fréquence par $2^{1/12}$ ; deux notes voisines
sont séparées si leur écart $f\,(2^{1/12}-1)$ dépasse $\Delta f$, donc au-dessus de
\begin{equation}
  f_{\min} = \frac{\Delta f}{2^{1/12}-1} \approx \num{16.8}\,\Delta f .
\end{equation}
Le nombre de trames par seconde vaut $\Fe/H = 2\Fe/N$. Le \cref{tab-resolutions} donne
ces grandeurs pour les quatre valeurs de $N$ demandées.

\begin{table}[h]
  \centering\small
  \caption{Résolutions de la TFCT à $\Fe = \qty{8}{\kilo\hertz}$, fenêtre de Hann
  ($\kw = \num{1.44}$), $H = N/2$. Dernière colonne : nombre $\Nt$ de valeurs entières de
  $\dm$ dans la zone cible $\Delta t \in [\num{0.1} ; 2]$\,s (\cref{sec-empreintes}).}
  \label{tab-resolutions}
  \begin{tabular}{@{}rrrrrr@{}}
    \toprule
    $N$ & $\Tw$ (ms) & $\Delta f$ (Hz) & trames/s & $f_{\min}$ (Hz) & $\Nt$ \\
    \midrule
    256  & 32  & 45          & \num{62.5}  & 757 & 119 \\
    512  & 64  & \num{22.5}  & \num{31.25} & 378 & 59 \\
    1024 & 128 & \num{11.25} & \num{15.6}  & 189 & 30 \\
    2048 & 256 & \num{5.6}   & \num{7.8}   & 95  & 15 \\
    \bottomrule
  \end{tabular}
\end{table}

\paragraph{Choix de $N = 1024$.} Quatre arguments vont dans le même sens.
\begin{enumerate}
  \item \emph{Fréquence.} Avec $\Delta f \approx \qty{11}{\hertz}$, deux demi-tons sont
        séparés dès \qty{190}{\hertz}, donc sur tout le registre mélodique. À $N = 512$, il
        faudrait monter à \qty{380}{\hertz} : les notes graves voisines se fondraient dans
        un même lobe et leurs pics deviendraient instables.
  \item \emph{Temps.} $\Tw = \qty{128}{\milli\second}$ reste plus court que la durée
        typique d'une note (\qtyrange{100}{500}{\milli\second}), si bien qu'une trame
        contient le plus souvent une seule configuration harmonique. À $N = 2048$, une
        trame de \qty{256}{\milli\second} mélangerait souvent deux notes successives.
  \item \emph{Spécificité des clés.} Le nombre de trames par seconde fixe le nombre $\Nt$
        de valeurs de $\dm$ possibles dans la zone cible, donc le nombre de clés $M$
        (\cref{sec-empreintes}). Passer à $N = 2048$ divise $\Nt$ et $M$ par deux.
  \item \emph{Robustesse au filtrage.} La propriété (2) exige que la réponse
        impulsionnelle soit courte devant la fenêtre (\cref{sec-proprietes}). Une fenêtre
        plus longue l'étend à des réponses plus longues : c'est le seul argument en faveur
        de $N$ plus grand, et il ne l'emporte pas sur les trois précédents.
\end{enumerate}
$N = 1024$ est le plus petit $N$ qui sépare les demi-tons sur tout le registre
mélodique ; tout $N$ plus grand dégrade la résolution temporelle et la spécificité sans
bénéfice musical.

\paragraph{Choix de $H = N/2$ et de $\Nfft = N$.} Avec $H = N/2$, les fenêtres de Hann
périodiques successives se recouvrent de moitié et leur somme est constante
($w[n] + w[n + N/2] = 1$) : chaque échantillon pèse autant dans l'analyse, et aucun
événement bref ne tombe dans un creux entre deux trames. Prendre $\Nfft > N$ (remplissage
par des zéros) interpolerait le spectre sans améliorer $\Delta f$, qui ne dépend que de $N$ :
on multiplierait les canaux, donc les clés possibles, sans information nouvelle, et les pics
passeraient plus facilement d'un canal à son voisin.

\subsection{Les deux propriétés qui fondent la méthode}
\label{sec-proprietes}

\paragraph{Propriété (1) : un décalage temporel translate le spectrogramme.}
Soit $y[n] = x[n - m_0H]$ avec $m_0$ entier. D'après (2),
\begin{equation}
\begin{split}
  X_y[m,k]
  &= \sum_{n=0}^{N-1} x[n + mH - m_0H]\, w[n]\, e^{-j2\pi kn/\Nfft} \\
  &= \sum_{n=0}^{N-1} x[n + (m-m_0)H]\, w[n]\, e^{-j2\pi kn/\Nfft}
  = X_x[m-m_0, k],
\end{split}
  \label{eq-translation}
\end{equation}
donc $S_y[m,k] = S_x[m-m_0,k]$ : le spectrogramme est translaté de $m_0$ trames. Le
résultat est exact, sans approximation, et vaut même pour la phase, parce que chaque
trame est analysée dans son propre repère temporel. C'est lui qui permet de retrouver la
position de l'extrait (\cref{sec-decision}).

\emph{Décalage quelconque.} Si $n_0 = m_0H + r$ avec $0 < r < H$, la trame $m$ de $y$
analyse un segment de $x$ qui commence $r$ échantillons avant la trame $m-m_0$ de sa
grille : elle tombe \emph{entre} deux trames de $x$, et~\eqref{eq-translation} ne
s'applique plus. Dans un passage stable, $|X|$ varie peu d'une position de fenêtre à
l'autre et les pics restent en place ; près d'une attaque ou d'une fin de note, le
maximum temporel d'un pic peut passer d'une trame à sa voisine. En rapportant chaque
trame de $y$ à la trame de $x$ la plus proche, l'écart est d'au plus $H/2 = N/4$
échantillons, soit un quart de fenêtre quel que soit $N$. Deux conséquences :
les votes du bon morceau se répartissent sur $\delta = m_0$ et $\delta = m_0 + 1$ au lieu
d'une seule case, ce qui abaisse le score ; et une paire dont un seul pic a migré voit son
$\dm$ changer d'une unité, donc sa clé est perdue.

\paragraph{Propriété (2) : un filtrage ajoute un biais qui ne dépend que de la fréquence.}
Soit $y = \alpha\,(h * x)$ avec $\alpha > 0$ et $h$ causal de longueur $L_h$. On prolonge
$w$ par zéro hors de $[0, N-1]$. En reportant $y[n+mH] = \alpha\sum_l h[l]\,x[n+mH-l]$
dans (2), puis en posant $n' = n - l$ :
\begin{align}
  X_y[m,k]
  &= \alpha \sum_{n} \sum_{l=0}^{L_h-1} h[l]\, x[n+mH-l]\, w[n]\, e^{-j2\pi kn/\Nfft} \nonumber\\
  &= \alpha \sum_{l=0}^{L_h-1} h[l]\, e^{-j2\pi kl/\Nfft}
     \sum_{n'} x[n'+mH]\, w[n'+l]\, e^{-j2\pi kn'/\Nfft} .
  \label{eq-filtrage-exact}
\end{align}
Si la fenêtre ne variait pas sur la durée de $h$, on pourrait remplacer $w[n'+l]$ par $w[n']$
et la somme intérieure deviendrait $X_x[m,k]$. On écrit donc $w[n'+l] = w[n'] + (w[n'+l] - w[n'])$ :
\begin{equation}
  X_y[m,k] = \alpha\,\hhat(f_k)\, X_x[m,k] + \varepsilon[m,k],
  \qquad
  \hhat(f_k) = \sum_{l=0}^{L_h-1} h[l]\, e^{-j2\pi f_k l/\Fe},
  \label{eq-filtrage}
\end{equation}
avec
\[
  \varepsilon[m,k] = \alpha \sum_{l} h[l]\, e^{-j2\pi kl/\Nfft}
  \sum_{n'} x[n'+mH]\,\bigl(w[n'+l] - w[n']\bigr)\, e^{-j2\pi kn'/\Nfft} .
\]
Le terme $\varepsilon$ vient de ce que la fenêtre ne voit pas exactement le même segment
de $x$ pour chaque retard $l$. Pour la fenêtre de Hann
$w[n] = \tfrac12\bigl(1 - \cos(2\pi n/N)\bigr)$, la pente est au plus $\pi/N$, et la fenêtre
est nulle à ses deux bords, donc continue une fois prolongée par zéro. Par conséquent
$|w[n'+l] - w[n']| \le \pi l/N$ pour tout $n'$, bords compris, et
\[
  |\varepsilon[m,k]| \;\le\; \alpha\,\frac{\pi}{N}
  \sum_{l} l\,|h[l]| \sum_{n'=-L_h}^{N-1} |x[n'+mH]| .
\]
Rapportée à $|X_x[m,k]|$, l'erreur est d'ordre $L_h/N$ : elle est négligeable si $L_h \ll N$.
(Avec une fenêtre rectangulaire, discontinue à ses bords, elle serait d'ordre 1 : la
continuité de la fenêtre de Hann est ici essentielle.) En prenant le module au carré puis le
logarithme, on obtient la relation (3) de l'énoncé :
\begin{equation}
  S_{y,\mathrm{dB}}[m,k] \approx S_{x,\mathrm{dB}}[m,k] + g(k),
  \qquad g(k) = 20\log_{10}\alpha + 20\log_{10}|\hhat(f_k)| .
  \tag{3}
\end{equation}
Le biais $g(k)$ ne dépend pas de la trame $m$ : dans le spectrogramme en décibels, le gain
et le filtrage de toute la chaîne se réduisent à un décalage vertical qui ne dépend que de
la fréquence. On ne garde que le module : la phase de $X_y$ contient un terme qui dépend
de $t_0$, inconnu, et la réverbération la rend aléatoire.

\paragraph{Vérification expérimentale (Q2).} Si (3) est vraie, $S_{y,\mathrm{dB}} - S_{x,\mathrm{dB}}$
tracé en fonction de $k$ donne la même courbe $g(k)$ pour toutes les trames, et cette
courbe suit $20\log_{10}|\hhat(f_k)|$. Pour un passe-bande de Butterworth
\qtyrange{300}{3400}{\hertz}, dont la réponse impulsionnelle décroît en quelques
millisecondes ($L_h \ll N = 1024$), les courbes doivent se superposer dans la bande
passante. Hors bande, $|X_y|$ tombe au niveau du bruit numérique et la différence devient
erratique : la relation ne vaut que là où le signal filtré reste au-dessus de ce plancher.
La \cref{fig-filtrage} le confirme : \aremplir{dispersion entre trames dans la bande
passante (en dB), écart à $20\log_{10}|\hhat|$, comportement hors bande}.

\begin{figure}[t]
  \centering
  \inclurefig[0.8\linewidth]{difference_spectres_passe_bande}{$S_{y,\mathrm{dB}} - S_{x,\mathrm{dB}}$ en fonction de k pour plusieurs trames, passe-bande 300-3400 Hz}
  \caption{Vérification de la propriété (2) : différence $S_{y,\mathrm{dB}} - S_{x,\mathrm{dB}}$
  en fonction de la fréquence pour \aremplir{5} trames d'un même morceau, après un filtre
  passe-bande \qtyrange{300}{3400}{\hertz}. En pointillés : $20\log_{10}|\hhat(f)|$.
  \aremplir{Ce qu'il faut y voir.}}
  \label{fig-filtrage}
\end{figure}

\subsection{Invariance des maxima locaux et condition (5)}
\label{sec-maxima}

La carte de constellation est l'ensemble des maxima locaux du spectrogramme (définition
(4), \cref{sec-pics}). Deux résultats justifient ce choix.

\paragraph{Invariance au gain.} Si $y = \alpha x$ avec $\alpha > 0$, alors
$X_y = \alpha X_x$ et $S_y = \alpha^2 S_x$. La fonction $s \mapsto \alpha^2 s$ est strictement
croissante, donc elle conserve l'ordre des valeurs : pour tout voisinage $V$,
\[
  S_y[m,k] = \max_{V(m,k)} S_y \iff S_x[m,k] = \max_{V(m,k)} S_x .
\]
L'ensemble des positions des maxima locaux est donc \emph{rigoureusement} identique pour
$x$ et $y$. On ne conserve que ces positions $(m,k)$, jamais les amplitudes, qui dépendent
de $\alpha$. Le seuil doit respecter la même invariance : un seuil local (médiane de $\SdB$
dans la bande, plus une marge) se décale de $20\log_{10}\alpha$ en même temps que les pics,
et un quota de pics par bande ne dépend que du classement des valeurs ; un seuil global
$\tau$ fixé en décibels ferait au contraire apparaître ou disparaître des pics avec le
volume. Seule la constante $\varepsilon$ de (2) échappe à l'invariance, mais elle
n'intervient que dans le silence numérique, où elle crée des plateaux de valeurs égales que
le test d'égalité prendrait pour des pics : le seuil local, strictement au-dessus de la
médiane, les élimine.

\paragraph{Invariance au filtrage : la condition (5).} Soit $(m,k_0)$ un pic de $S_x$ et
\[
  \mu = S_{x,\mathrm{dB}}[m,k_0] - \max_{(m',k') \in V^*} S_{x,\mathrm{dB}}[m',k'] \;>\; 0,
  \qquad V^* = V(m,k_0) \setminus \{(m,k_0)\},
\]
son avance en décibels sur son concurrent le plus fort dans le voisinage. Notons
\[
  \Delta g(V) = \max_{|k'-k_0| \le \kv} |g(k') - g(k_0)|
\]
la variation du biais sur la largeur du voisinage. D'après (3), pour tout
$(m',k') \in V^*$,
\[
  S_{y,\mathrm{dB}}[m,k_0] - S_{y,\mathrm{dB}}[m',k']
  = \underbrace{S_{x,\mathrm{dB}}[m,k_0] - S_{x,\mathrm{dB}}[m',k']}_{\ge\, \mu}
  + \underbrace{g(k_0) - g(k')}_{\ge\, -\Delta g(V)}
  \;\ge\; \mu - \Delta g(V) .
\]
Le pic reste donc le maximum de son voisinage dès que
\begin{equation}
  \Delta g(V) < \mu \quad \Longrightarrow \quad \text{le pic survit au filtrage.}
  \tag{5}
\end{equation}
La condition est suffisante, pas nécessaire : un pic peut survivre avec
$\Delta g(V) \ge \mu$ si son concurrent le plus proche en valeur n'est pas celui que le filtre
favorise. Elle est toujours remplie pour un gain seul ($g$ constant, $\Delta g = 0$). Pour un
filtre dont la réponse varie de $|g'|$ dB par canal, $\Delta g(V) \approx |g'|\,\kv$, qui ne
dépend que de la \emph{demi-largeur fréquentielle} $\kv$ du voisinage : la largeur
temporelle $\mv$ n'intervient pas, puisque $g$ ne dépend pas de $m$.

On en tire la règle de dimensionnement de l'énoncé. Un voisinage étroit en fréquence
donne un $\Delta g$ petit, donc des pics robustes au filtrage, mais il produit beaucoup de
pics, dont certains peu dominants ($\mu$ faible). Un voisinage large ne garde que des pics
très dominants, donc une empreinte compacte et spécifique, mais il exige
$|g'|\,\kv < \mu$ sur une plus grande largeur. Dans la bande passante du filtre
\qtyrange{300}{3400}{\hertz}, $g$ est presque plat et presque tous les pics survivent ; près
des fréquences de coupure, la pente atteint plusieurs décibels sur $\pm 5$ canaux, et seuls
les pics nettement dominants restent en place.

\subsection{Limites des deux propriétés}
\label{sec-limites}

\paragraph{Réverbération.} La propriété (2) suppose $L_h \ll N$. Une réverbération de temps
$T_{60}$ décroît de \qty{60}{\decibel} en $T_{60}$, donc de \qty{20}{\decibel} en $T_{60}/3$ ;
sa longueur utile est $L_h \approx T_{60}\Fe/3$, soit environ 800 échantillons pour
$T_{60} = \qty{0.3}{\second}$ (comparable à $N$ : cas limite) et \num{2700} pour
$T_{60} = \qty{1}{\second}$ (plus du double de $N$). La borne sur $\varepsilon$ n'est plus
petite, et la courbe $S_{y,\mathrm{dB}} - S_{x,\mathrm{dB}}$ dépend de la trame. Une meilleure
image est celle d'un \emph{bruit additif produit par le signal lui-même} : la queue
réverbérée d'une note se superpose aux trames suivantes. Elle comble les creux entre les
notes, efface les pics des notes brèves et laisse survivre ceux des notes tenues, plus
énergétiques que la traîne.

\paragraph{Décalage non entier.} La propriété (1) n'est exacte que pour un décalage multiple
de $H$ ; dans le cas réaliste, une partie des pics migre d'une trame
(\cref{sec-proprietes}).

\paragraph{Tempo et hauteur.} Ni (1) ni (2) ne couvrent un changement de tempo, qui
modifie les écarts $\dm$, ni un changement de hauteur, qui multiplie toutes les fréquences
par un même facteur et change donc les canaux $k$ et les écarts $\dk$ : la méthode n'y est
pas invariante.

% !TEX root = ../main.tex
% ---------------------------------------------------------------------
\section{Conception}
\label{sec-conception}
% ---------------------------------------------------------------------

Chaque étape est justifiée par les propriétés de la \cref{sec-theorie} ; les valeurs finales
des paramètres sont rassemblées au \cref{tab-parametres}.

\subsection{Carte de constellation (Q3)}
\label{sec-pics}

\paragraph{Définition.} Pour un voisinage $V(m,k)$ de $(2\mv+1)$ trames sur $(2\kv+1)$
canaux centré en $(m,k)$, la carte de constellation est
\begin{equation}
  \mathcal{P} = \Bigl\{ (m,k) : S[m,k] = \max_{V(m,k)} S \ \text{ et } \ \SdB[m,k] \ge \tau[m,k] \Bigr\},
  \tag{4}
\end{equation}
où $\tau[m,k]$ est un seuil \emph{local}. D'après la \cref{sec-maxima}, cette définition est
exactement invariante au gain, et invariante au filtrage sous la condition (5).

\paragraph{Voisinage de départ.} Avec $\mv = 3$ et $\kv = 5$, un pic domine
$\pm 3 \times \qty{64}{\milli\second} = \pm\qty{192}{\milli\second}$ et
$\pm 5 \times \qty{7.8}{\hertz} = \pm\qty{39}{\hertz}$. La demi-largeur en fréquence vaut
un demi-ton vers \qty{650}{\hertz} : elle est assez étroite pour que $\Delta g(V)$
reste petit, et assez large pour qu'un même lobe spectral (largeur
$\Delta f \approx \qty{11}{\hertz}$, soit 1,4 canal) ne donne qu'un seul pic.

\paragraph{Seuil local et quota par bande.} Le seuil $\tau$ est la médiane de $\SdB$ dans la
bande de fréquence, augmentée d'une marge : il suit le gain et le niveau moyen de la bande,
donc il est invariant au gain et quasi invariant à un filtrage lent. On impose en plus un
nombre fixe de pics par bande et par seconde, en gardant les plus forts : la densité $\rho$
est ainsi maîtrisée, quels que soient le volume et le genre. Les 6 bandes sont en
progression géométrique, $f_b = f_{\min}(f_{\max}/f_{\min})^{b/N_b}$ entre 200 et
\qty{4000}{\hertz}, de rapport constant $20^{1/6} \approx \num{1.65}$ (environ $\num{0.72}$
octave, soit 8,6 demi-tons par bande) : chaque bande couvre le même intervalle musical. Un
découpage linéaire placerait la moitié des pics imposés par le quota au-dessus de
\qty{2.1}{\kilo\hertz}, là où l'énergie musicale est faible et les pics fragiles.

\paragraph{Question 3 : pourquoi $\rho p^2$ et pas $p$ seul.} Notons $\rho$ la densité de
pics (pics/s) et $p$ la probabilité qu'un pic survive à la dégradation. Une requête de
durée $\Tq$ contient environ $\rho\Tq$ ancres ; chacune forme $F$ paires, et une paire
n'est retrouvée que si ses deux pics survivent. Si les survies des deux pics sont
indépendantes, une clé survit avec la probabilité $p^2$ et le nombre de clés correctes
retrouvées vaut
\begin{equation}
  \E[\text{clés correctes}] = F\,\rho\,\Tq\,p^2 .
  \label{eq-cles-correctes}
\end{equation}
À $F$ et $\Tq$ fixés, c'est donc $\rho p^2$ qui mesure ce qu'un réglage du voisinage apporte
au score du bon morceau. Le critère $p$ seul est trompeur, car $p$ et $\rho$ varient en sens
inverse : un voisinage \emph{large} ne garde que des pics très dominants, qui survivent bien
($p$ grand) mais sont rares ($\rho$ petit) ; un voisinage \emph{étroit} donne beaucoup de pics,
dont une bonne partie, faiblement dominants, sont déplacés par le bruit ($p$ petit).
Maximiser $p$ conduirait à une empreinte presque vide. Le carré compte aussi : une clé exige
la survie de \emph{deux} pics, donc une baisse de $p$ pèse plus qu'une hausse équivalente de
$\rho$. L'hypothèse d'indépendance est d'ailleurs plutôt pessimiste : deux pics proches
subissent le même bruit local, et une clé survit alors avec une probabilité supérieure à
$p^2$.

\paragraph{Mesure.} Avec le voisinage de départ ($\mv = 3$, $\kv = 5$), le quota de
30~pics/s et une marge de \qty{6}{\decibel}, on obtient $\rho \approx 25$~pics/s : le quota
n'est pas atteint dans les passages où peu de maxima dépassent le seuil local. Nous avons balayé le voisinage sans quota, de $\pm 1$ à $\pm 5$ trames
et de $\pm 1$ à $\pm 12$ canaux, sur 100 extraits de \qty{8}{\second} (un par morceau)
bruités à \qty{10}{\decibel} (\cref{fig-p-rho}). Comme prévu, $\rho$ et $p$ varient en sens
inverse : quand le voisinage grandit, $\rho$ chute de 308 à 23~pics/s, et $p$ ne monte que
de \num{0.30} à \num{0.42}. $\rho p^2$ décroît donc de \num{27.7} à \num{4.1}. Le réglage
retenu ($\pm 3$ trames $\times$ $\pm 5$ canaux, avec quota) donne $\rho = 26$~pics/s,
$p = \num{0.55}$ et $\rho p^2 = \num{7.8}$, soit près de deux fois plus que le voisinage
sans quota de même densité ($p = \num{0.42}$, $\rho p^2 = \num{4.1}$) : le quota garde les
pics les plus forts de chaque bande, qui sont ceux qui dépassent le plus le bruit.
Maximiser $\rho p^2$ pousserait vers des voisinages minuscules, mais la taille de l'index
croît comme $\rho F$ et les appariements fortuits comme $\rho^2$ : on fixe donc $\rho$ dans
la plage du sujet (25 à 35~pics/s) par le quota, et c'est à densité fixée que le critère
$\rho p^2$ départage les réglages.

\begin{figure}[t]
  \centering
  \inclurefig[0.8\linewidth]{p_en_fonction_de_rho}{p sous bruit à 10 dB en fonction de rho, un point par voisinage}
  \caption{Proportion $p$ de pics conservés après ajout d'un bruit blanc à
  \qty{10}{\decibel}. Voisinages de $\pm 1$, $\pm 3$ et $\pm 5$ trames sur 3 à 25 canaux,
  sans quota ; étoile : réglage retenu, avec quota. 100 extraits de \qty{8}{\second}.
  (a)~$\rho$ en fonction de la largeur du voisinage ; (b)~$p$ en fonction de $\rho$ ;
  (c)~critère $\rho p^2$. $p$ augmente quand le voisinage grandit, mais beaucoup moins vite
  que $\rho$ ne diminue ; à densité égale, le quota double $\rho p^2$.}
  \label{fig-p-rho}
\end{figure}

\subsection{Empreintes par paires de pics (Q4)}
\label{sec-empreintes}

\paragraph{La clé.} À chaque pic ancre $(m_1,k_1)$ on associe les $F$ premiers pics cibles
$(m_2,k_2)$ de la zone $\dm_{\min} \le m_2 - m_1 \le \dm_{\max}$, $|k_2 - k_1| \le \dk_{\max}$,
et chaque paire donne la clé
\begin{equation}
  u = \bigl(Q_f(k_1),\, Q_f(k_2),\, \dm\bigr), \qquad \dm = m_2 - m_1, \qquad Q_f(k) = \lfloor k/\gf \rfloor ,
  \tag{6}
\end{equation}
rangée dans l'index avec la valeur (morceau, $m_1$). La clé ne dépend que de positions de
pics, donc elle est invariante au gain ; elle ne contient que l'écart $m_2 - m_1$, qu'une
translation (propriété (1)) laisse inchangé, et c'est pour cela que l'instant absolu $m_1$
est rangé dans la valeur et non dans la clé. Le triplet est codé en un seul entier
$u = (Q_f(k_1)\,\Nf + Q_f(k_2))\,\Nt + (\dm - \dm_{\min})$, inférieur à
$M = \Nf^2\Nt \le 2^{32}$.

\paragraph{Zone cible.} À \num{15.6} trames/s, $\Delta t \in [\num{0.1} ; 2]$\,s donne
$\dm \in [2, 31]$, soit $\Nt = 30$ valeurs. On exclut $\dm = 0$ (pics d'un même accord, très
fréquents, peu spécifiques) et $\dm = 1$ (où la migration d'un pic d'une trame change le plus
souvent la clé). La borne $\dk_{\max} = 60$ canaux (\qty{470}{\hertz}) garde la cible dans la
même région musicale. La zone contient en moyenne
$30 \times \num{1.9} \times 121/490 \approx 14$ pics ; $F$ doit rester en dessous pour que les
ancres des bords de bande aient assez de cibles.

\paragraph{Pourquoi des paires.} Une clé « 1 pic » survit avec la probabilité $p$, une clé
« 1 paire » avec $p^2$, mais il y a $F$ fois plus de paires que de pics : le nombre de clés
correctes est multiplié par $Fp$, et les paires gagnent dès que $F\,p > 1$. Dans le même
temps, le nombre de clés possibles passe de $K = 513$ à $M = \Nf^2 \Nt \approx \num{7.9e6}$ :
les paires sont plus nombreuses à survivre \emph{et} \num{15000} fois plus discriminantes, au
prix d'un index $F$ fois plus gros (\cref{tab-pic-paire}).

\begin{table}[h]
  \centering\small
  \caption{Clés « 1 pic » contre clés « 1 paire », mesurées sur 100 requêtes de
  $\Tq = \qty{8}{\second}$ alignées, $\rho = \num{26.1}$~pics/s, $F = 10$. Colonne de gauche :
  exemple de l'énoncé ($\Tq = \qty{8}{\second}$, $\rho = 30$, $F = 10$, $p = \num{0.35}$).
  Nombre de paires mesuré : \num{1707}, et non $F\rho\Tq = \num{2090}$, car les ancres de fin
  d'extrait n'ont pas toutes leurs $F$ cibles.}
  \label{tab-pic-paire}
  \begin{tabular}{@{}lccccc@{}}
    \toprule
    & & \multicolumn{2}{c}{Bruit blanc \qty{10}{\decibel}} & \multicolumn{2}{c}{MP3 \qty{64}{\kilo\bit\per\second}} \\
    \cmidrule(lr){3-4}\cmidrule(l){5-6}
    & Énoncé (pic / paire) & 1 pic & 1 paire & 1 pic & 1 paire \\
    \midrule
    Clés dans la requête
      & 240 / 2400 & 209 & \num{1707} & 209 & \num{1707} \\
    Probabilité de survie : $p$ / $p^2$
      & \num{0.35} / \num{0.12} & \num{0.56} & \num{0.31} & \num{0.86} & \num{0.73} \\
    Survie des clés mesurée
      & -- & -- & \num{0.36} & -- & \num{0.75} \\
    Clés correctes retrouvées (mesurées)
      & 84 / 294 & 116 & 612 & 179 & \num{1276} \\
    Clés possibles : $K$ / $M$
      & 513 / \num{7.9e6} & 513 & \num{7.9e6} & 513 & \num{7.9e6} \\
    \bottomrule
  \end{tabular}
\end{table}

La condition $Fp > 1$ est largement remplie : $Fp = \num{5.6}$ sous bruit à
\qty{10}{\decibel} et \num{8.6} en MP3. Les paires donnent 4 à 7 fois plus de clés correctes
que les pics isolés (612 contre 116, \num{1276} contre 179), tout en étant \num{15000} fois plus
discriminantes. La survie des clés mesurée dépasse un peu $p^2$ (\num{0.36} contre
\num{0.31} sous bruit) : les deux pics d'une paire ne disparaissent pas indépendamment, les
pics forts survivent ensemble et les pics faibles disparaissent ensemble. L'hypothèse
d'indépendance est donc légèrement pessimiste, comme annoncé à la \cref{sec-pics}.

\paragraph{Regroupement en fréquence : $\gf$.} Regrouper $\gf$ canaux voisins en une classe
donne $\Nf = \lceil K/\gf \rceil$ classes et $M = \Nf^2\Nt$ clés possibles. Si les clés étaient
réparties uniformément, une requête de $\nq$ clés interrogée contre un index de $\Nu$ entrées
produirait en moyenne $\nq\Nu/M$ appariements fortuits (\cref{sec-decision}). Le
\cref{tab-gammaf} donne ces grandeurs pour $\gf = 1, 2, 4, 8$.

\begin{table}[h]
  \centering\small
  \caption{Effet du regroupement $\gf$. Parasites attendus : modèle uniforme $\nq\Nu/M$, avec
  $\nq \approx \num{1015}$ clés mesurées par requête de \qty{5}{\second} et
  $\Nu = \num{677171}$ (notre index de 100 extraits de \qty{30}{\second}). Plancher : seuil
  prévu par le modèle uniforme (\cref{eq-plancher}). $\sthr$ mesuré : seuil recalibré pour
  chaque $\gf$ sur 50 négatives (plage selon le RSB). \TBD{} : $\Tq = \qty{5}{\second}$,
  200 requêtes positives ; $\TFD = 0/50$ partout.}
  \label{tab-gammaf}
  \begin{tabular}{@{}rrrrrrrr@{}}
    \toprule
    $\gf$ & $\Nf$ & $M$ & Parasites & Plancher & $\sthr$ mesuré & \TBD{} \qty{10}{\decibel} & \TBD{} \qty{0}{\decibel} \\
    \midrule
    1 & 513 & \num{7.9e6} & 87         & 3 & 21 à 26 & \qty{95.5}{\percent} & \qty{81.5}{\percent} \\
    2 & 257 & \num{2.0e6} & 347        & 3 & 26 à 33 & \qty{96.0}{\percent} & \qty{81.5}{\percent} \\
    4 & 129 & \num{5.0e5} & \num{1380} & 4 & 35 à 41 & \qty{96.0}{\percent} & \qty{80.5}{\percent} \\
    8 & 65  & \num{1.3e5} & \num{5430} & 6 & 59      & \qty{94.0}{\percent} & \qty{71.0}{\percent} \\
    \bottomrule
  \end{tabular}
\end{table}

\paragraph{Question 4 : le meilleur $\gf$ n'est pas celui qui donne le moins de parasites.}
Les parasites croissent comme $\gf^2$, donc $\gf = 1$ les minimise. Mais deux effets jouent
en sens inverse. D'abord, les parasites gênent peu tant qu'ils restent dispersés : la
décision porte sur la hauteur du plus grand pic de l'histogramme des $\delta$, pas sur le
nombre total d'appariements, et le maximum de nombreuses variables de Poisson de petite
moyenne croît très lentement (le plancher de seuil ne passe que de 3 à 6 quand les
parasites sont multipliés par 64). Ensuite, un regroupement absorbe les petits déplacements
de pics : si un pic survivant se décale d'un canal avec la probabilité $j$, il ne change de
classe qu'avec la probabilité $j/\gf$, et une clé survit avec la probabilité
$\bigl[p\,(1 - j/\gf)\bigr]^2$. Pour $j = \num{0.3}$, passer de $\gf = 1$ à 2 multiplie les clés
correctes par environ \num{1.5}. Le bon critère est donc la \emph{séparation} entre le score
du bon morceau et celui des mauvais. Quand $\gf$ devient grand, une même clé se répète dans
chaque morceau, les votes fortuits s'empilent dans les mêmes cases et la spécificité
s'effondre : l'optimum est intermédiaire.

\paragraph{Mesure.} Le \cref{tab-gammaf} confirme les deux effets, mais le second est bien
plus fort que prévu : le score médian du bon morceau à \qty{0}{\decibel} augmente avec $\gf$
(49, 65, 76, 88), mais le seuil nécessaire pour garder $\TFD = 0$ monte de 21 à 59, et non de
3 à 6, car les clés réelles sont loin d'être uniformes (\cref{sec-decision}). D'où un
plateau : $\gf = 1$ et 2 donnent le même \TBD{}, puis le \TBD{} à \qty{0}{\decibel} chute à
\qty{71}{\percent} pour $\gf = 8$ (à \qty{10}{\decibel}, les écarts ne valent que 1 à 4
requêtes sur 200). Nous gardons $\gf = 1$ : même \TBD{} que $\gf = 2$, quatre fois moins
d'appariements fortuits et le seuil le plus bas.

\paragraph{Choix de $F$.} Avec les mêmes requêtes, $F = 5$, 10 et 15 donnent
\qtylist{78.0;81.5;81.5}{\percent} de \TBD{} à \qty{0}{\decibel}, pour un index de
\numlist{368190;677171;905862} entrées. Le score augmente avec $F$, mais le seuil aussi :
le \TBD{} plafonne à partir de $F = 10$, alors que l'index continue de croître. $F = 10$ est
le plus petit $F$ qui atteint ce plafond.

\subsection{Index inversé}
\label{sec-index}

L'index inversé range, pour chaque clé, la liste des couples (morceau, $m_1$) : une requête
ne lit que ses $\nq$ clés au lieu de se comparer à chaque morceau. Nous le stockons comme un
tableau $[u,\ \text{morceau},\ m_1]$ trié par $u$ ; chaque clé est cherchée par dichotomie,
puis \emph{toutes} ses entrées sont lues. Le coût d'une requête est
$O(\nq \log \Nu + A)$, où seul le nombre $A$ d'appariements lus dépend de la taille de la
base. La taille attendue, $\Nu = N_b\,L\,\rho\,F \approx \num{7.5e5}$ entrées pour 100
extraits de \qty{30}{\second}, est proche des \num{677171} mesurées (environ 225 par seconde
de musique, les ancres de fin d'extrait ayant moins de $F$ cibles), soit
\qty{14.6}{\mega\byte}. La \cref{fig-passage-echelle} confirme les deux prédictions : la taille
croît linéairement (de \qty{1.8}{\mega\byte} pour 10 morceaux à \qty{56}{\mega\byte} pour 500)
et le temps de requête médian reste presque constant (\qtyrange{1.9}{2.4}{\milli\second}). Ces temps comptent l'identification à partir du signal prétraité, pas le décodage du
fichier.

\begin{figure}[t]
  \centering
  \inclurefig[0.7\linewidth]{temps_et_taille_index}{temps de requête et taille de l'index en fonction du nombre de morceaux}
  \caption{Passage à l'échelle, de 10 à 500 morceaux indexés. (a)~Temps médian d'une
  requête de \qty{5}{\second} non dégradée : quasi constant. (b)~Taille de l'index en
  mémoire : proportionnelle au nombre de morceaux (tirets : \qty{0.11}{Mo} par morceau).}
  \label{fig-passage-echelle}
\end{figure}

\subsection{Alignement temporel et décision (Q5)}
\label{sec-decision}

\paragraph{Alignement.} Si la requête commence à la trame $m_0$ de la référence, la
propriété (1) assure qu'une paire trouvée dans la requête à l'ancre $m_q$ l'est dans la
référence en $m_r = m_q + m_0$. Tous les appariements corrects vérifient donc
\begin{equation}
  \delta = m_r - m_q = m_0 ,
  \tag{7}
\end{equation}
alors que les appariements fortuits donnent des $\delta$ dispersés. Le score $s$ d'un
morceau est le maximum de son histogramme des $\delta$ ; on répond « le morceau de score
maximal » si $s \ge \sthr$, « inconnu » sinon. Sur la chaîne « jouet »
(\cref{fig-histo-delta}), 699 appariements du bon morceau tombent dans une seule case, en
$\delta$ = \qty{6.66}{\second} pour une position réelle de \qty{6.65}{\second}, alors qu'un
mauvais morceau n'en réunit jamais plus de 2. Les pics secondaires du bon morceau viennent
vraisemblablement d'une boucle musicale répétée.

\begin{figure}[t]
  \centering
  \inclurefig[0.75\linewidth]{histogramme_delta}{histogramme des delta : mauvais morceau contre bon morceau}
  \caption{Histogramme des décalages $\delta = m_r - m_q$ pour une requête de
  \qty{8}{\second} non dégradée (chaîne « jouet »). À gauche, un mauvais morceau : votes
  dispersés, score 2. À droite, le bon morceau : pic net de 699 votes en $\delta = m_0$.}
  \label{fig-histo-delta}
\end{figure}

\paragraph{Score du bon morceau.} Toutes les paires correctes votent dans la même case,
donc d'après~\eqref{eq-cles-correctes}
\begin{equation}
  s_{\text{bon}} \approx F\,\rho\,\Tq\,p^2 ,
  \label{eq-score-bon}
\end{equation}
soit environ 300 pour l'exemple de l'énoncé et $1250\,p^2$ avec nos valeurs
($\rho = 25$, $\Tq = \qty{5}{\second}$). C'est une borne haute : les ancres de fin de requête
n'ont pas toutes leurs cibles ($\nq \approx \num{1015}$ paires et non \num{1250}), et même sans
dégradation $p < 1$, car l'extrait ne tombe pas sur la grille des trames. La
\cref{sec-theorie-mesure} mesure ces deux effets.

\paragraph{Score d'un mauvais morceau.} Si les clés étaient uniformes, les appariements
fortuits avec un mauvais morceau suivraient une loi de Poisson de moyenne
$\nq N_c/M \approx \num{1.1}$ (pour $N_c \approx \num{6800}$ entrées par extrait), répartie
sur environ 550 cases de $\delta$ : un seuil de 3 suffirait à garder la probabilité de fausse
détection sous \qty{1}{\percent} (\cref{eq-plancher}, \cref{sec-annexe-scores}). C'est un
\emph{plancher}. Nous mesurons pourtant des scores négatifs jusqu'à 26 (de 10 à 25 d'après
l'énoncé), car le modèle uniforme suppose deux choses fausses. D'abord, \emph{les clés ne
sont pas uniformes} : la gamme tempérée concentre les pics sur les notes, et les intervalles
consonants reviennent dans tous les morceaux ; il faut remplacer $M$ par un alphabet effectif
$\Meff = 1/\sum_u \pi_u^2$ bien plus petit. Ensuite, \emph{les appariements fortuits ne sont
pas indépendants} : une note tenue produit la même clé à plusieurs ancres consécutives, et une
note tenue commune à la requête et à un mauvais morceau donne plusieurs votes dans la
\emph{même} case, exactement comme un bon morceau. Deux ou trois coïncidences de ce type
suffisent à expliquer des scores de 10 à 25. Le seuil doit donc être calibré sur des
requêtes négatives.

\paragraph{Question 5 : fixer $\sthr$ automatiquement.} On cherche le plus petit seuil qui
garantit un \TFD{} visé $\alpha$ (critère de Neyman-Pearson). La méthode :
(1)~séparer les négatives en un jeu de calibration et un jeu de test, extraits de morceaux
différents ; (2)~ajuster sur les scores de calibration une loi des valeurs extrêmes, la loi
de Gumbel~\cite{gumbel1958}, par la méthode des moments ; (3)~l'inverser pour obtenir le
seuil qui donne $\Prob(S > \sthr) = \alpha$, ce qui extrapole au-delà du plus grand score
observé (\cref{sec-annexe-scores}) ; (4)~valider sur les négatives de test. Cette validation
est limitée par le nombre $n$ de négatives : si aucune n'est acceptée, la borne à
\qty{95}{\percent} du \TFD{} vaut environ $3/n$, soit \qty{6}{\percent} pour $n = 50$ ; il
faudrait 300 négatives sans erreur pour garantir moins de \qty{1}{\percent}. Le modèle indique
aussi que le seuil ne croît qu'avec le logarithme du nombre de morceaux.

\paragraph{Mise en œuvre et résultat.} Le plus fort score des 50 négatives de calibration,
toutes conditions confondues, vaut 23 ; nous avons retenu
$\sthr = \lceil \num{1.5} \times 23 \rceil = 35$. La marge est nécessaire : sur les négatives de
\emph{test}, le score monte jusqu'à 26, ce que seul le découpage calibration/test révèle. Avec
$\sthr = 35$, $\TFD = 0/50$ dans toutes les conditions. La loi de Gumbel, testée, donne
$\sthr \approx 16$ à 18 pour $\alpha = \qty{1}{\percent}$, mais un \TFD{} réel de 4 à
\qty{5}{\percent} (\cref{tab-gumbel}) : sa queue est trop légère, pour les deux raisons
ci-dessus. Elle donne l'ordre de grandeur ; le seuil prudent de 35 garantit l'absence de
fausse détection.

\paragraph{Séparation des scores.} La \cref{fig-histo-scores} montre que les deux
populations sont séparées d'un ordre de grandeur : à $\Tq = \qty{5}{\second}$, les négatives
ne dépassent pas 23, les positives vont pour la plupart de 30 à 900. Les queues se
recouvrent entre 10 et 25 : quelques positives très dégradées (bruit à \qty{0}{\decibel},
réverbération à \qty{1}{\second}) deviennent des non-détections, jamais des confusions. Sur
la \cref{fig-tbd-tfd}, ces deux conditions sont les seules où il faut arbitrer : un seuil
choisi pour elles seules (18 à 20) donnerait 83 et \qty{82}{\percent} de \TBD{} à $\TFD = 0$,
contre \num{66.5} et \qty{51}{\percent} avec le seuil unique, qui doit tenir pour toutes les
conditions à la fois.

\begin{figure}[t]
  \centering
  \inclurefig[0.78\linewidth]{histogrammes_scores}{histogrammes des scores positifs et négatifs, abscisse log, seuil}
  \caption{Histogrammes des scores des requêtes positives et négatives, avec le seuil
  $\sthr = 35$. $\Tq = \qty{5}{\second}$, toutes dégradations réunies (brouhaha compris) :
  \num{3000} requêtes positives et \num{1500} négatives. (a)~Vue d'ensemble en abscisse
  logarithmique ; les scores étant entiers, les classes sont bornées sur des demi-entiers et
  ramenées à une même largeur en log. (b)~Zoom linéaire de 0 à 60 : aucune négative entre
  23 et 35, c'est la marge du seuil.}
  \label{fig-histo-scores}
\end{figure}

\begin{figure}[t]
  \centering
  \inclurefig[0.72\linewidth]{courbe_tbd_tfd}{TBD en fonction du TFD en balayant le seuil, une courbe par dégradation}
  \caption{\TBD{} en fonction du \TFD{}, obtenu en balayant le seuil, pour six conditions,
  $\Tq = \qty{5}{\second}$, $N^+ = 200$. La courbe ne dépend pas du seuil retenu : elle
  utilise les 100 négatives (calibration et test), d'où un \TFD{} qui avance par pas de
  \qty{1}{\percent}. (a) Vue d'ensemble ; (b) zoom sur le coin supérieur gauche, où les courbes
  des conditions faciles se confondent avec le bord. Les cercles marquent le seuil retenu,
  $\sthr = 35$.}
  \label{fig-tbd-tfd}
\end{figure}

\subsection{Récapitulatif des paramètres}
\label{sec-parametres}

Le \cref{tab-parametres} rassemble la valeur retenue pour chaque paramètre et la mesure ou
la propriété qui la justifie.

\begin{table}[!htb]
  \centering\footnotesize
  \caption{Paramètres utilisés pour toutes les mesures, et leur justification. Toutes les
  valeurs sont celles de départ du sujet ; les balayages confirment qu'aucune autre valeur
  testée ne fait mieux.}
  \label{tab-parametres}
  \begin{tabularx}{\linewidth}{@{}llX@{}}
    \toprule
    Paramètre & Valeur & Justification \\
    \midrule
    $\Fe$ & \qty{8}{\kilo\hertz} & Shannon, bande utile \qtyrange{0}{4}{\kilo\hertz} (\cref{sec-shannon}) \\
    Fenêtre & Hann périodique & lobes secondaires faibles ; continue, donc (3) vaut avec une erreur d'ordre $L_h/N$ (\cref{sec-proprietes}) \\
    $N$ & 1024 & sépare les demi-tons dès \qty{190}{\hertz}, $\Tw = \qty{128}{\milli\second}$ (\cref{sec-tfct}) \\
    $H$, $\Nfft$ & $N/2$, $N$ & somme des fenêtres constante ; pas d'interpolation inutile (\cref{sec-tfct}) \\
    $\mv$, $\kv$ & 3, 5 & condition (5) ; avec le quota, $\rho p^2 = \num{7.8}$, le double d'un voisinage sans quota de même $\rho$ (\cref{fig-p-rho}) \\
    Seuil, quota & médiane + \qty{6}{\decibel}, 5 pics/bande/s & invariance au gain, maîtrise de $\rho$ ; $\rho$ mesuré $\approx 25$~pics/s (\cref{sec-pics}) \\
    Bandes & 6, géométriques, \qtyrange{200}{4000}{\hertz} & même intervalle musical par bande (\cref{sec-pics}) \\
    Zone cible & $\dm \in [2, 31]$, $|\dk| \le 60$ & exclut les paires instables ; $\Nt = 30$ (\cref{sec-empreintes}) \\
    $F$ & 10 & $Fp > 1$ ; plus petit $F$ qui atteint le plateau de \TBD{} (5, 10, 15 testés, \cref{sec-empreintes}) \\
    $\gf$ & 1 & même \TBD{} que $\gf = 2$, moins de parasites ; $\gf = 8$ fait chuter le \TBD{} (\cref{tab-gammaf}) \\
    $\Tq$ & \qty{5}{\second} & score $\propto \Tq$ (\cref{eq-score-bon}) ; 3 et \qty{10}{\second} aussi mesurés \\
    $\sthr$ & 35 & $\lceil \num{1.5} \times{}$ plus fort score négatif de calibration$\rceil$ ; $\TFD = 0/50$ en test (\cref{sec-decision}) \\
    \bottomrule
  \end{tabularx}
\end{table}

% ---------------------------------------------------------------------
\section{Résultats}
\label{sec:resultats}
% ---------------------------------------------------------------------
% Squelette (partie B3, à rédiger à partir de resultats/notes_resultats.md).
% Règle : chaque taux est donné avec son seuil, sa durée Tq, sa dégradation et son
% nombre de requêtes.

\subsection{Protocole}
\label{sec:protocole}

\aremplir{Base : 100 morceaux de \texttt{fma\_small} répartis sur 8 genres ; requêtes
positives (nombre, tirage de l'instant de début) et négatives (nombre, morceaux hors base) ;
ordre extrait $\to$ dégradation $\to$ prétraitement ; calcul de $P_x$ sur l'extrait (8) ;
séparation calibration / test des négatives ; graines aléatoires.}

La définition (8) du rapport signal à bruit,
\begin{equation}
  \RSB_{\mathrm{dB}} = 10\log_{10}\frac{P_x}{P_b},
  \tag{8}
\end{equation}
est appliquée avec $P_x$ calculé sur l'extrait servant de requête, et non sur le morceau
entier : sinon, un extrait calme recevrait un bruit relatif bien plus fort que le \RSB{}
annoncé.

\begin{figure}[t]
  \centering
  \inclurefig[0.9\linewidth]{constellation}{spectrogramme avec la constellation superposée}
  \caption{Spectrogramme d'un extrait de \aremplir{morceau} avec la carte de constellation
  superposée.}
  \label{fig:constellation}
\end{figure}

\subsection{Taux de bonne et de fausse détection}
\label{sec:tbd-tfd}

\aremplir{TBD et TFD par dégradation, distinction « mauvais morceau » / « inconnu ».}

\begin{figure}[t]
  \centering
  \inclurefig[0.8\linewidth]{tbd_en_fonction_du_rsb}{TBD en fonction du RSB, une courbe par Tq = 3, 5, 10 s}
  \caption{\TBD{} en fonction du \RSB{} (bruit blanc), une courbe par durée
  $\Tq \in \{3, 5, 10\}$\,s, $\sthr = \aremplir{}$.}
  \label{fig:tbd-rsb}
\end{figure}

\begin{figure}[t]
  \centering
  \inclurefig[0.7\linewidth]{tbd_en_fonction_du_debit_mp3}{TBD en fonction du débit MP3}
  \caption{\TBD{} en fonction du débit MP3, $\Tq = \aremplir{}$\,s, $\sthr = \aremplir{}$.}
  \label{fig:tbd-mp3}
\end{figure}

\subsection{Conditions réelles}
\label{sec:reel}

\aremplir{Enregistrements au micro à deux distances, brouhaha : protocole (salle, matériel,
distances) et résultats.}

\subsection{Passage à l'échelle}

\begin{figure}[t]
  \centering
  \inclurefig[0.8\linewidth]{temps_et_taille_index}{temps de requête et taille de l'index en fonction du nombre de morceaux}
  \caption{Temps de requête et taille de l'index en fonction du nombre de morceaux indexés.}
  \label{fig:passage-echelle}
\end{figure}

\aremplir{Comparer à la \cref{sec:index} : taille proportionnelle à $N_b$, temps de requête
presque constant.}

% ---------------------------------------------------------------------
\section{Discussion des cas d'échec}
\label{sec:discussion}
% ---------------------------------------------------------------------

\subsection{\texorpdfstring{Pourquoi le MP3 à \qty{64}{\kilo\bit\per\second} gêne moins qu'un bruit blanc à \qty{10}{\decibel} (Q6)}{Pourquoi le MP3 à 64 kbit/s gêne moins qu'un bruit blanc à 10 dB (Q6)}}
\label{sec:q6}

\aremplir{Constat mesuré : \TBD{} = x~\% en MP3 \qty{64}{\kilo\bit\per\second} contre y~\% sous
bruit blanc à \qty{10}{\decibel} ($\Tq$, $\sthr$, nombre de requêtes), et $p$ = x contre y
(\cref{tab:pic-paire}).} Deux dégradations de « niveau » comparable ont donc des effets très
différents. L'explication tient à la manière dont chacune répartit son bruit dans le plan
temps-fréquence, comparée à la définition (4) des pics.

\paragraph{Le codage perceptuel épargne les pics.} Un codeur MP3 découpe le signal en
sous-bandes, estime pour chaque trame un \emph{seuil de masquage}, c'est-à-dire le niveau
en dessous duquel l'oreille ne perçoit rien en présence des composantes fortes voisines, et
alloue ses bits de façon à maintenir le bruit de quantification sous ce
seuil~\cite{painter2000}. Les composantes situées sous le seuil sont grossièrement
quantifiées ou supprimées. Or ce seuil est construit à partir des \emph{masqueurs},
c'est-à-dire des composantes les plus fortes de leur voisinage temps-fréquence ; il se
situe typiquement de \qtyrange{10}{20}{\decibel} sous un masqueur tonal. Un pic au sens de
(4) est, par définition, le maximum de son voisinage $V$ : c'est précisément un masqueur.
Le codeur le conserve, et le bruit qu'il ajoute autour reste nettement plus bas que lui.
Dans les termes de la \cref{sec:maxima}, l'avance $\mu$ du pic sur ses voisins n'est que
peu entamée, et il reste en place. Le codage supprime surtout les composantes faibles,
qui ne sont pas des pics : il simplifie le spectrogramme sans toucher à sa constellation.
Enfin, le bruit de quantification suit l'enveloppe spectrale du signal : le rapport signal
à bruit \emph{local} est à peu près le même dans toutes les bandes.

\paragraph{Le bruit blanc noie les pics des bandes hautes.} Un bruit blanc de variance
$\sigma_b^2$ ajoute en chaque point du spectrogramme une contribution de puissance moyenne
\[
  \E\bigl|B[m,k]\bigr|^2 = \sigma_b^2 \sum_{n=0}^{N-1} w[n]^2 ,
\]
\emph{la même dans tous les canaux}. Le spectre musical, lui, décroît avec la fréquence :
l'essentiel de l'énergie est sous \qty{1}{\kilo\hertz}. Le \RSB{} de \qty{10}{\decibel} est
un rapport \emph{global}, dominé par les basses fréquences ; le rapport local dans le canal
$k$ vaut
\[
  \RSB_{\text{local}}(k) = \RSB + 10\log_{10} \frac{S_x(k)}{\overline{S_x}} ,
\]
où $\overline{S_x}$ est la moyenne du spectre sur tous les canaux. Dans une bande où le
spectre musical est \qty{20}{\decibel} sous sa moyenne, ce qui est courant entre 2 et
\qty{4}{\kilo\hertz}, le rapport local tombe à \qty{-10}{\decibel} : le bruit y domine, et
ses propres fluctuations créent des maxima locaux qui remplacent ceux de la musique. Le
quota par bande aggrave l'effet : il impose de prendre le même nombre de pics dans les
bandes hautes que dans les basses, et ce sont alors des pics de bruit, différents à chaque
requête. Une bonne partie des pics des deux ou trois bandes supérieures est perdue, donc
$p$ chute, et le score, qui varie en $p^2$, chute encore plus.

\paragraph{En résumé.} Le MP3 est conçu pour préserver ce qui domine localement, et c'est
exactement ce que retiennent les pics. Le bruit blanc ne tient aucun compte du signal, et
son effet est maximal là où la musique est faible. \aremplir{Si la mesure de $p$ par bande
est disponible, l'ajouter ici : $p$ à peu près constant selon la bande en MP3, décroissant
vers les hautes fréquences sous bruit blanc.} À \qty{32}{\kilo\bit\per\second}, le codeur
manque de bits : le bruit de quantification dépasse le seuil de masquage, les composantes
moyennes disparaissent par blocs et certains pics secondaires avec elles, ce qui explique
la baisse du \TBD{} à bas débit (\cref{fig:tbd-mp3}).

\subsection{Genres difficiles}
\label{sec:genres}

\aremplir{Constat mesuré : \TBD{} par genre (\cref{fig:genres}), genres les plus et les
moins bien reconnus, avec $\Tq$, $\sthr$ et la dégradation.}

La théorie permet d'anticiper quels genres posent problème, et pourquoi.

\begin{itemize}
  \item \emph{Musiques tonales et tenues (Folk, Instrumental, une partie d'International).}
        Les notes sont stables, leurs harmoniques dépassent nettement le fond : les pics ont
        une forte avance $\mu$ et survivent bien au bruit et au filtrage. On attend un bon
        \TBD{}. En contrepartie, les notes tenues produisent des clés répétées sur plusieurs
        trames, donc des votes fortuits cohérents (\cref{sec:decision}) : les scores
        négatifs peuvent y être plus élevés.
  \item \emph{Musiques à spectre dense ou bruité (Rock saturé, Experimental).} Une guitare
        saturée ou une texture bruitée remplit le spectre de façon presque uniforme : la
        platitude spectrale est élevée, les maxima locaux dépassent à peine leurs voisins
        ($\mu$ faible), et le moindre bruit les déplace. La condition (5) et son équivalent
        pour le bruit sont rarement remplies : $p$ est faible.
  \item \emph{Musiques dominées par les percussions et la voix parlée (Hip-Hop).} L'énergie
        est concentrée dans des transitoires brefs et dans les basses ; une voix parlée a
        une hauteur qui glisse en permanence. Les pics sont peu nombreux dans les bandes
        moyennes et hautes, ou instables en fréquence, et un décalage non aligné sur la
        grille (\cref{sec:proprietes}) déplace facilement un transitoire d'une trame.
  \item \emph{Musiques en boucles (Electronic, Hip-Hop).} Une boucle répétée produit les
        mêmes clés à chaque répétition : l'histogramme des $\delta$ du bon morceau présente
        plusieurs pics, ce qui ne gêne pas l'identification mais rend la position estimée
        ambiguë. Des sons de synthèse ou des échantillons de batterie partagés entre
        morceaux peuvent en revanche créer des coïncidences cohérentes avec de mauvais
        morceaux.
\end{itemize}
\aremplir{Confronter ces hypothèses à la mesure : les genres prévus comme difficiles le
sont-ils ? Exemple d'échec analysé : spectrogramme et constellation d'une requête ratée.}

\begin{figure}[t]
  \centering
  \inclurefig[0.8\linewidth]{tbd_par_genre}{TBD par genre}
  \caption{\TBD{} par genre pour \aremplir{dégradation, $\Tq$, $\sthr$, nombre de requêtes
  par genre}.}
  \label{fig:genres}
\end{figure}

\subsection{Réverbération et enregistrements réels}
\label{sec:echecs-reverb}

La réverbération sort du cadre de la propriété (2) dès que sa réponse dépasse la durée de
la fenêtre (\cref{sec:limites}). Elle agit comme un bruit additif fabriqué par le signal
lui-même : on attend une dégradation modérée à $T_{60} = \qty{0.3}{\second}$ et nette à
$T_{60} = \qty{1}{\second}$, plus marquée pour les morceaux rapides que pour les morceaux
lents. \aremplir{Résultats de la réverbération et des enregistrements en salle à 1 et
4~m : \TBD{} mesurés, et comparaison avec les dégradations simulées.}

\subsection{Limites de la méthode}

La méthode n'est invariante ni au changement de tempo (les écarts $\dm$ changent) ni au
changement de hauteur (les canaux changent) : une version accélérée, transposée ou jouée
en direct n'est pas reconnue. Elle identifie un \emph{enregistrement}, pas une
\emph{œuvre}.

% ---------------------------------------------------------------------
\section{Conclusion}
\label{sec:conclusion}
% ---------------------------------------------------------------------
% Squelette (partie B3).

\aremplir{Bilan : ce que le système fait, avec ses chiffres clés ; ce que la théorie
prévoyait (invariance au gain et au filtrage, score en $F\rho\Tq p^2$, plancher de seuil,
supériorité du MP3 sur le bruit blanc) et ce que la mesure a confirmé ou contredit.}

\aremplir{Perspectives : empreintes invariantes au tempo et à la hauteur (rapports de
fréquences et rapports d'écarts temporels au lieu de différences), comparaison avec
l'approche par signes de différences d'énergie de Haitsma et Kalker, base plus grande.}


\nocite{muller2015}
\printbibliography[heading=bibintoc,title={Références}]

\clearpage
\appendix
% !TEX root = ../main.tex
% ---------------------------------------------------------------------
\section{Démonstrations}
\label{sec-demos}
% ---------------------------------------------------------------------

\subsection{Propriété (1) : un décalage entier translate le spectrogramme}
\label{sec-demo-translation}

Soit $y[n] = x[n - m_0H]$ avec $m_0$ entier. D'après (2),
\begin{equation}
\begin{split}
  X_y[m,k]
  &= \sum_{n=0}^{N-1} x[n + mH - m_0H]\, w[n]\, e^{-j2\pi kn/\Nfft} \\
  &= \sum_{n=0}^{N-1} x[n + (m-m_0)H]\, w[n]\, e^{-j2\pi kn/\Nfft}
  = X_x[m-m_0, k],
\end{split}
  \label{eq-translation}
\end{equation}
donc $S_y[m,k] = S_x[m-m_0,k]$. Le résultat est exact et vaut même pour la phase, parce que
chaque trame est analysée dans son propre repère temporel. Si $n_0 = m_0H + r$ avec
$0 < r < H$, la trame $m$ de $y$ analyse un segment de $x$ qui commence $r$ échantillons
avant la trame $m-m_0$ : en rapportant chaque trame de $y$ à la trame de $x$ la plus proche,
l'écart est d'au plus $H/2 = N/4$ échantillons, et~\eqref{eq-translation} n'est plus
qu'approchée.

\subsection{Propriété (2) : un filtrage court ajoute un biais \texorpdfstring{$g(k)$}{g(k)}}
\label{sec-demo-filtrage}

Soit $y = \alpha\,(h * x)$ avec $\alpha > 0$ et $h$ causal de longueur $L_h$. On prolonge
$w$ par zéro hors de $[0, N-1]$. En reportant $y[n+mH] = \alpha\sum_l h[l]\,x[n+mH-l]$
dans (2), puis en posant $n' = n - l$ :
\begin{equation}
  X_y[m,k]
  = \alpha \sum_{l=0}^{L_h-1} h[l]\, e^{-j2\pi kl/\Nfft}
     \sum_{n'} x[n'+mH]\, w[n'+l]\, e^{-j2\pi kn'/\Nfft} .
  \label{eq-filtrage-exact}
\end{equation}
Si la fenêtre ne variait pas sur la durée de $h$, on pourrait remplacer $w[n'+l]$ par $w[n']$
et la somme intérieure deviendrait $X_x[m,k]$. En écrivant
$w[n'+l] = w[n'] + (w[n'+l] - w[n'])$, on obtient
$X_y[m,k] = \alpha\,\hhat(f_k)\, X_x[m,k] + \varepsilon[m,k]$ avec
\[
  \varepsilon[m,k] = \alpha \sum_{l} h[l]\, e^{-j2\pi kl/\Nfft}
  \sum_{n'} x[n'+mH]\,\bigl(w[n'+l] - w[n']\bigr)\, e^{-j2\pi kn'/\Nfft} .
\]
Pour la fenêtre de Hann $w[n] = \tfrac12\bigl(1 - \cos(2\pi n/N)\bigr)$, la pente est au
plus $\pi/N$, et la fenêtre est nulle à ses deux bords, donc continue une fois prolongée par
zéro. Par conséquent $|w[n'+l] - w[n']| \le \pi l/N$ pour tout $n'$, bords compris, et
\[
  |\varepsilon[m,k]| \;\le\; \alpha\,\frac{\pi}{N}
  \sum_{l} l\,|h[l]| \sum_{n'=-L_h}^{N-1} |x[n'+mH]| .
\]
Rapportée à $|X_x[m,k]|$, l'erreur est d'ordre $L_h/N$. Avec une fenêtre rectangulaire,
discontinue à ses bords, elle serait d'ordre 1 : la continuité de la fenêtre de Hann est ici
essentielle. Le module au carré puis le logarithme donnent la relation (3).

\subsection{Invariance des maxima au gain et condition (5)}
\label{sec-demo-maxima}

\paragraph{Gain.} Si $y = \alpha x$, alors $S_y = \alpha^2 S_x$, et $s \mapsto \alpha^2 s$ est
strictement croissante : pour tout voisinage $V$,
$S_y[m,k] = \max_{V(m,k)} S_y \iff S_x[m,k] = \max_{V(m,k)} S_x$.

\paragraph{Filtrage.} Soit $(m,k_0)$ un pic de $S_x$, d'avance $\mu$ sur son concurrent le
plus fort dans $V^* = V(m,k_0) \setminus \{(m,k_0)\}$. D'après (3), pour tout
$(m',k') \in V^*$,
\[
  S_{y,\mathrm{dB}}[m,k_0] - S_{y,\mathrm{dB}}[m',k']
  = \underbrace{S_{x,\mathrm{dB}}[m,k_0] - S_{x,\mathrm{dB}}[m',k']}_{\ge\, \mu}
  + \underbrace{g(k_0) - g(k')}_{\ge\, -\Delta g(V)}
  \;\ge\; \mu - \Delta g(V) ,
\]
qui est strictement positif sous la condition (5) : le pic reste le maximum de son voisinage.
La condition est suffisante, pas nécessaire : un pic peut survivre avec $\Delta g(V) \ge \mu$
si son concurrent le plus proche en valeur n'est pas celui que le filtre favorise.

\subsection{Score d'un mauvais morceau et seuil de Gumbel}
\label{sec-annexe-scores}

\paragraph{Modèle uniforme.} Si les clés étaient réparties uniformément sur les $M$ clés
possibles, chacune des $\nq$ clés de la requête trouverait en moyenne $N_c/M$ entrées d'un
morceau $c$ qui en compte $N_c$. Le nombre d'appariements fortuits avec ce morceau suit
alors une loi de Poisson de moyenne $\lambda_c = \nq N_c/M$, répartie sur les
$D \approx (L + \Tq)\Fe/H$ cases de son histogramme des $\delta$ (environ 550 pour des
extraits de \qty{30}{\second}). Chaque case reçoit une loi de Poisson de moyenne
$\lambda = \lambda_c/D$, et
\begin{equation}
  \Prob(\text{fausse détection}) \;\le\; (N_b - 1)\, D\; \Prob\bigl(\mathrm{Poisson}(\lambda) \ge s\bigr) .
  \label{eq-plancher}
\end{equation}
Avec $\nq = 1250$, $N_c \approx \num{6800}$ et $\gf = 1$, $\lambda_c \approx \num{1.1}$, et
$s = 3$ suffit à rester sous \qty{1}{\percent}.

\paragraph{Alphabet effectif.} Si $\pi_u$ est la proportion des entrées de l'index qui
portent la clé $u$, et si les clés de la requête suivent la même loi, la probabilité qu'une
clé de la requête tombe sur une entrée donnée vaut $\sum_u \pi_u^2$ au lieu de $1/M$. Le
modèle reste valable en remplaçant $M$ par $\Meff = 1/\sum_u \pi_u^2 \le M$, avec égalité
seulement si les clés sont uniformes. La contrainte $|\dk| \le 60$ et les bandes
\qtyrange{200}{4000}{\hertz} réduisent déjà l'alphabet accessible à environ \num{1.8e6} clés,
quatre fois moins que $M$.

\paragraph{Notes tenues.} Si la requête contient la clé $u$ à $L$ ancres consécutives
$m, \dots, m+L-1$ et un mauvais morceau à $L'$ ancres consécutives $r, \dots, r+L'-1$, les
$LL'$ appariements donnent $\delta = (r - m) + (i - j)$, et la case $\delta = r - m$ en reçoit
$\min(L, L')$ : des votes cohérents dans le temps, comme pour un bon morceau.

\paragraph{Seuil de Gumbel.} La loi de Gumbel a pour fonction de répartition
$\Prob(S \le s) = \exp\bigl(-e^{-(s-\mu_G)/\beta}\bigr)$. Par la méthode des moments,
\[
  \beta = \frac{\sqrt{6}}{\pi}\,\sigma_s , \qquad \mu_G = \bar{s} - \gamma_E\,\beta ,
  \qquad \gamma_E \approx \num{0.5772},
\]
où $\bar{s}$ et $\sigma_s$ sont la moyenne et l'écart-type des scores de calibration. Pour un
\TFD{} visé $\alpha$, $\Prob(S > \sthr) = \alpha$ donne
$\sthr = \mu_G - \beta \ln\bigl(-\ln(1 - \alpha)\bigr) \approx \mu_G + \num{4.6}\,\beta$ pour
$\alpha = \qty{1}{\percent}$. Passer de $N_b$ à $N_b'$ morceaux décale la loi de
$\beta\ln(N_b'/N_b)$. Mesures (\cref{tab-gumbel}) : le \TFD{} réel est 4 à 5 fois plus élevé
que le \TFD{} visé.

\begin{table}[h]
  \centering\small
  \caption{Seuil de Gumbel et \TFD{} réel mesuré sur les négatives de test des 24 conditions
  (\num{1200} scores).}
  \label{tab-gumbel}
  \begin{tabular}{@{}lcccccc@{}}
    \toprule
    Calibration & $\bar{s}$ & $\sigma_s$ & $\sthr$ ($\alpha = \qty{1}{\percent}$) & \TFD{} réel & $\sthr$ ($\alpha = \qty{0.1}{\percent}$) & \TFD{} réel \\
    \midrule
    50 scores, sans dégradation & \num{7.86} & \num{2.71} & \num{16.4} & \qty{5.0}{\percent} & \num{21.2} & \qty{1.6}{\percent} \\
    \num{1200} scores, toutes conditions & \num{8.09} & \num{3.01} & \num{17.5} & \qty{4.3}{\percent} & \num{22.9} & \qty{1.1}{\percent} \\
    \bottomrule
  \end{tabular}
\end{table}

% ---------------------------------------------------------------------
\section{Notations}
\label{sec-notations}
% ---------------------------------------------------------------------

Les notations sont celles de l'énoncé ; celles que nous ajoutons sont signalées par un
astérisque.

\begin{center}\small
\begin{tabularx}{\linewidth}{@{}lX@{}}
  \toprule
  Symbole & Signification \\
  \midrule
  $x(t)$, $y(t)$ & morceau d'origine, extrait capté \\
  $\alpha$, $h$, $t_0$, $b(t)$ & gain, réponse impulsionnelle haut-parleur $\to$ salle $\to$ micro, décalage, bruit additif (1) \\
  $\Fe$ & fréquence d'échantillonnage \\
  $N$, $H$, $\Nfft$ & longueur de fenêtre, avance entre trames, points de la TFD \\
  $K = \Nfft/2 + 1$ & nombre de canaux \\
  $S[m,k]$, $\SdB[m,k]$ & spectrogramme (trame $m$, canal $k$), en décibels (2) \\
  $t_m = mH/\Fe$, $f_k = k\Fe/\Nfft$ & instant de la trame $m$, fréquence du canal $k$ \\
  $\Tw$, $\Delta f$, $\kw$ & résolution temporelle, résolution fréquentielle, constante de la fenêtre \\
  $g(k)$ & biais introduit par le gain et le filtrage (3) \\
  $V$, $\mv$, $\kv$ & voisinage de $(2\mv+1)$ trames sur $(2\kv+1)$ canaux \\
  $\mu$, $\Delta g(V)$ & avance d'un pic sur ses voisins, variation de $g$ sur le voisinage (5) \\
  $\rho$ & densité de pics (pics/s) \\
  $p$ & probabilité qu'un pic survive à la dégradation \\
  $F$ & nombre de cibles par ancre \\
  $u$, $Q_f$, $\gf$ & clé, classe de fréquence, nombre de canaux par classe (6) \\
  $\dm$, $\dk$ & écarts en trames et en canaux entre ancre et cible \\
  $M = \Nf^2\Nt$ & nombre de clés possibles, $\Nf$ classes de fréquence, $\Nt$ valeurs de $\dm$ \\
  $\Tq$ & durée de la requête \\
  $\nq$, $\Nu$ & nombre de clés de la requête, nombre d'entrées de l'index \\
  $\delta = m_r - m_q$ & décalage entre instants d'ancre dans la référence et dans la requête (7) \\
  $s$, $\sthr$ & score (maximum de l'histogramme des $\delta$), seuil de décision \\
  $N^+$, $N^-$ & nombres de requêtes positives et négatives \\
  \TBD{}, \TFD{} & taux de bonne détection, taux de fausse détection \\
  $\RSB_{\mathrm{dB}}$, $P_x$, $P_b$ & rapport signal à bruit, puissances du signal et du bruit (8) \\
  $L_h$* & longueur de la réponse impulsionnelle $h$ \\
  $N_b$*, $D$* & nombre de morceaux de la base, nombre de cases de l'histogramme des $\delta$ \\
  $\Meff$* & alphabet effectif $1/\sum_u \pi_u^2$ \\
  $\mu_G$*, $\beta$* & paramètres de la loi de Gumbel des scores négatifs \\
  \bottomrule
\end{tabularx}
\end{center}

\section{Liste des morceaux}
\label{sec-liste-morceaux}

Les 200 morceaux utilisés, tous extraits de \texttt{fma\_small}~\cite{fma2017} (extraits de
\qty{30}{\second}, licences Creative Commons). Le fichier complet, avec la licence et la durée
de chaque morceau, est \texttt{liste\_morceaux.csv}. Les titres trop longs sont abrégés, et
quelques caractères non latins ont été translittérés.

% Généré automatiquement depuis liste_morceaux.csv : ne pas modifier à la main.
\begin{footnotesize}
\begin{longtable}{@{}rlll@{}}
\toprule Identifiant & Artiste & Titre & Genre \\ \midrule \endhead
\bottomrule \endfoot
\multicolumn{4}{@{}l}{\textbf{Morceaux de la base (100)}} \\ \midrule
26652 & Graphiqs Groove & Bleuacide & Electronic \\
48439 & Bass Assassin & Panic & Electronic \\
55808 & Blasterhead & Killbots & Electronic \\
67638 & Southman & Wake Up & Electronic \\
69732 & Godblesscomputers & Natürlich & Electronic \\
80793 & Ostin & Spirit & Electronic \\
84198 & Lee Rosevere & Illuminations & Electronic \\
84484 & BLEO & Tart (Pts 1 \& 2) feat. KeFF & Electronic \\
105674 & Gilo & Outro & Electronic \\
107427 & UncleBibby & Brutish & Electronic \\
108415 & UncleBibby & AstroTurf & Electronic \\
142079 & Mike B. Fort & OOOTHERSIDE - Crystal Air & Electronic \\
143046 & Mystery Mammal & Apollo & Electronic \\
22150 & Jeff McLeod & The Horror Approaches & Experimental \\
29752 & Mister Jan & Row Row Row Your Boat (Mister Jan) & Experimental \\
34994 & 6 or 7 & Dear Listener & Experimental \\
39278 & U Can Unlearn Guitar & Chicago, Etc. & Experimental \\
48861 & NaDa BaBa & Skull Flower & Experimental \\
49064 & Justice Yeldham & Track 06 & Experimental \\
52861 & Ilya Monosov & For Jojo Takayanagi & Experimental \\
58054 & Trance And The Arcade & Without Resistance: Nature's Ni\ldots & Experimental \\
68539 & Benboncan & Foyer Crowd & Experimental \\
74392 & Thicker Than Water & Joy to the World & Experimental \\
116514 & SULCUS & R.A.H.M.A.V.N. [221] & Experimental \\
120782 & Teleg. & things are not as they seem & Experimental \\
126521 & L o W \& CROCODILE TEARS & Hymne de apophia & Experimental \\
10443 & Big Blood & Beast (1 Beast) & Folk \\
10993 & Big Blood & Oh Country (Skin \&amp; Bones) & Folk \\
29971 & Amy Sayer & wild tape & Folk \\
47663 & Onda bidón & me aborrezco cada vez que llego\ldots & Folk \\
67016 & Michael Chapman & Kodak Ghosts & Folk \\
88870 & Josh Woodward & Methylchloroisothiazolinone (No\ldots & Folk \\
91443 & Cal Folger Day & Easy Sometimes & Folk \\
108299 & Carter Behnke & Intro & Folk \\
116873 & Itasca & After Dawn & Folk \\
137907 & Matús Novanský & Sally & Folk \\
138416 & Derek Clegg & Egypt, Maybe Rome & Folk \\
139110 & Anthony Pasquarosa & What Makes A Man? & Folk \\
148211 & Josh Armistead & Summertime & Folk \\
11770 & 6th Sense & Future Music Instrumental     \{\ldots & Hip-Hop \\
11775 & 6th Sense & We Don't Stop Instrumental\ldots & Hip-Hop \\
11861 & 6th Sense \&amp; Mick Boogie & Nike Boots Freestyle & Hip-Hop \\
13768 & Majik Most & Just a Feelin' & Hip-Hop \\
14542 & Opio & Due Damage & Hip-Hop \\
19438 & Just Plain Ant & Stumble & Hip-Hop \\
26014 & Comfort Fit & Whenever You Wanna Call Me & Hip-Hop \\
30043 & Cullah & Soul Hip Hop & Hip-Hop \\
72135 & C-Doc & Life Slow Down (Instrumental) & Hip-Hop \\
97037 & CuzOH! Black Feat. Wordsmith & Stronger Living Feat Wordsmith & Hip-Hop \\
108838 & Tha Silent Partner & On Harry Love Drums & Hip-Hop \\
149689 & K. Sparks & Purpose & Hip-Hop \\
150288 & The Impossebulls & Thank U ('92) & Hip-Hop \\
63291 & I Was A Teenage Strangler Sound\ldots & March of the Teenagers pt. 2 & Instrumental \\
86487 & Andrea Carri & Sunshine After Rain & Instrumental \\
92953 & Muhmood & Organ Etude & Instrumental \\
119898 & Blue Dot Sessions & Ferry Landing & Instrumental \\
120298 & Cory Gray & Attic & Instrumental \\
120315 & Cory Gray & Life Size Head & Instrumental \\
125659 & Blue Dot Sessions & Demalion - Theme & Instrumental \\
125681 & Blue Dot Sessions & Eventual Victory & Instrumental \\
125820 & Blue Dot Sessions & Master & Instrumental \\
129042 & Blue Dot Sessions & Third in Line & Instrumental \\
133942 & Raul Diaz Palomar & Plan II & Instrumental \\
141282 & Kevin Hartnell & Rogue Planet & Instrumental \\
6360 & Orkestra Keyif & Sabah Sabah Seyredelim Yaliyi & International \\
6366 & Zikrayat & Taheya & International \\
36272 & Demiran Cerimovic and His Orkestar & Competition Rehearsal, Song 1 & International \\
36273 & Demiran Cerimovic and His Orkestar & Competition Rehearsal, Song 2 & International \\
41018 & Unity Vibration & Au royaume de la money & International \\
54468 & DuBoLoGy & Dub Action & International \\
59718 & Black Sea Hotel & Spavamise & International \\
60865 & Gogofski & Medley: Prishtinka, Ramo-Ramo,\ldots & International \\
62589 & Tchakare Kanyembe & Tchakare Kanyembe track 01 & International \\
73761 & Dad Rocks! & Weapons & International \\
114940 & Zingeray! & Di fishelekh in vaser & International \\
117942 & Volfoniq & L'amour est enfant dub & International \\
22480 & Ergo Phizmiz & Celia & Pop \\
34511 & Symfoniorkestern & Vilse (Slutet Cover) & Pop \\
42911 & Phemale & MEDICAL SEX & Pop \\
70425 & Supercute & Assked & Pop \\
90592 & Pilesar & Bullets Have No Effect & Pop \\
120105 & Lanark & Happy Boccherini & Pop \\
121740 & Frozen Houses & finished enconomic survivor & Pop \\
131452 & Mom Jeans & B06 Say We're Sweethearts Again & Pop \\
131454 & Mom Jeans & B08 Mario & Pop \\
135339 & Caleb Lemond & I Think From Now On & Pop \\
148286 & LoveyDove & Beppe of Italy & Pop \\
148288 & LoveyDove & Deep Down Inc. & Pop \\
42851 & Waylon Thornton and the Heavy H\ldots & Manson Halloween & Rock \\
44918 & Apache Dropout & Teenager & Rock \\
55718 & et\_ & Kopeika & Rock \\
68821 & Zulus & Death in the Current & Rock \\
69781 & Home Blitz & Secret Wave & Rock \\
96400 & Insane Ride & Save Your Tears & Rock \\
110634 & Eidetic Seeing & No Pilot & Rock \\
119257 & The Wades & like love & Rock \\
122397 & Gorilla Chief & Cocaine Girl & Rock \\
125001 & Let Me Crazy & 899 Seconds & Rock \\
126300 & Andy G. Cohen & Scramby Eggs & Rock \\
142511 & Miët & Ratatatatouille & Rock \\
\midrule
\multicolumn{4}{@{}l}{\textbf{Morceaux hors base, requêtes négatives (100)}} \\ \midrule
23016 & The Lights Galaxia & While She Sleeps (Morning Edit) & Electronic \\
48465 & Sycamore Drive & Ocean Breeze & Electronic \\
52409 & Buskerdroid & I Want You & Electronic \\
52630 & Jason Shaw & Quirky Jerk & Electronic \\
65063 & Sycamore Drive & Kicks & Electronic \\
66390 & krackatoa & Komodo Dodo & Electronic \\
69211 & JMD v Kube & Interference (Psychadelik Pedes\ldots & Electronic \\
72067 & Decktonic & Africa & Electronic \\
85831 & Frame & Flames (The Beginning) & Electronic \\
87070 & (F.E.V.E.R.) & Argonaut Into Deep Space (SFI r\ldots & Electronic \\
92366 & Out Of Orion (Ox3) & Horror House & Electronic \\
126899 & SPCZ & Kuma & Electronic \\
146458 & Langax & The Cavern & Electronic \\
38891 & File Under Toner & justified & Experimental \\
44781 & Logan Whitehurst and The Jr. Sc\ldots & (She's) Ten Feet Tall & Experimental \\
49068 & Justice Yeldham & Knobranaut & Experimental \\
51158 & Base Cleft & Carousel of Bent Facts & Experimental \\
56015 & Tone Poets & Surface & Experimental \\
63161 & Lucky Dragons & 13 & Experimental \\
73821 & Miquel Parera Jaques & nxVacuity002-2012-12-12\_04\_00\_16 & Experimental \\
110611 & Jeff Kolar & Happy Birthday To You & Experimental \\
112668 & ODOR BABY & 45 & Experimental \\
112788 & ODOR BABY & BROKEN NEWS & Experimental \\
124870 & Ross Oldenburg & Five-O (Improvisation for Five\ldots & Experimental \\
133294 & AG Davis & Jazz Got Fuck'd & Experimental \\
145043 & Hermit & Rock'n'Roll Hara-Kiri & Experimental \\
12058 & Big Blood & Land Upon Mud I & Folk \\
12060 & Big Blood & Heart of Glass & Folk \\
12064 & Big Blood & Beatle Bones 'N' Smokin' Stones & Folk \\
13930 & Foghorn Stringband & John Henry & Folk \\
41192 & Weyes Blood & Storms That Breed & Folk \\
63253 & Big Blood & Sirens Knell & Folk \\
69765 & Plusplus & Echoes & Folk \\
89477 & Wellington Sea Shanty Society & Great Open Sea & Folk \\
130955 & Albin Andersson & Even If It's True & Folk \\
135044 & Anga & Le diable & Folk \\
140791 & NaDa BaBa & Sam's Song & Folk \\
144475 & Nick Jaina & The Dinner Party (Instrumental) & Folk \\
145551 & This Mess is Mine & Thief & Folk \\
22091 & ETHX & Wondering Why (Featuring Nate M\ldots & Hip-Hop \\
43867 & Chicken Jones & one last try & Hip-Hop \\
56498 & Opio & World of Pain & Hip-Hop \\
63755 & Cullah & A Dangerous Game & Hip-Hop \\
63803 & Lushlife & Gymnopedie (feat. Shad) & Hip-Hop \\
71690 & C-Doc & Drumm Machine (12 Inch Mixx) (F\ldots & Hip-Hop \\
75372 & Kellee Maize & 2012 & Hip-Hop \\
75425 & Kellee Maize & Eleven 11 Remix & Hip-Hop \\
97697 & Tickle & Stand Up & Hip-Hop \\
116390 & Asthmatic Astronaut & Jazz soul & Hip-Hop \\
129914 & Dollar Boyz & King Of Da City & Hip-Hop \\
132962 & Cockmasta Flex & Maricones Calientes (remastered) & Hip-Hop \\
140623 & The Impossebulls & The More Things Change & Hip-Hop \\
54153 & Mark Neil & 05 TheLastDodo TheHappyAndroid & Instrumental \\
75435 & Cherly KaCherly & Time is running out. I'm sure. & Instrumental \\
114212 & Squire Tuck & Old Fashioned Lullaby & Instrumental \\
114534 & Blue Dot Sessions & Diner Gothic - Dasher & Instrumental \\
116177 & Squire Tuck & A Mexican Love Affair & Instrumental \\
121651 & happiness in aeroplanes & Her History of Water & Instrumental \\
123761 & Hawkin & Week 1 - rain & Instrumental \\
123764 & Hawkin & Week 4 - clarity & Instrumental \\
126411 & Blue Dot Sessions & Chromium Blush & Instrumental \\
142641 & TRG Banks & Goodbye Machineryhead & Instrumental \\
144454 & TRG Banks & Father Green & Instrumental \\
148132 & Ryan Cullinane & Upside of Looking Down Instrume\ldots & Instrumental \\
1688 & The Sounds of Taraab & Sishituki & International \\
44822 & Black Bear Combo & Watchyerstep & International \\
54469 & DuBoLoGy & Walk In Forest & International \\
54482 & Friglob & Love music & International \\
60753 & Niva & Snosti ti Dojdov & International \\
80753 & Gogofski & Medley: Sadila Jana, Davaj Me & International \\
99362 & Svitanya & Vet Me The & International \\
108867 & Black Sea Hotel & Malka Moma Dvora Mela & International \\
129926 & Radio Jarocho and Zenen Zeferin\ldots & Vieja & International \\
131941 & Choban Elektrik & Unknown Cocek Tune & International \\
150080 & DubRaJah & Deep Karnataka & International \\
152253 & Panda Dub & Milky Way & International \\
28485 & The Pop Winds & Good Person & Pop \\
32334 & Plushgoolash & Brown Rainbows & Pop \\
37369 & Útidúr & Grasping for Thoughts & Pop \\
59445 & Phemale & This Place & Pop \\
61670 & Gosprom & Japan Punk & Pop \\
61736 & Soft Serve & Margaritaville & Pop \\
68353 & Caroline Gorman & Hiding in the Dark & Pop \\
91160 & Monroeville Music Center & Invaders From Within & Pop \\
127184 & Wintermitts & I Saw A Ghost & Pop \\
129391 & A la Mode & Mountaintop & Pop \\
134385 & Scott Holmes & Moving Forward & Pop \\
138038 & Massage & Breaking Up & Pop \\
13556 & STEEP & Let Me Go & Rock \\
69723 & ROGER ROBERT & Dracula & Rock \\
93986 & Bugs and Rats & Irish Clouds & Rock \\
106277 & Night Rally & A Birthday Party & Rock \\
106570 & Apache Tomcat & Texas Ostrich Farm & Rock \\
108502 & Voodoo Puppets & Slender Man & Rock \\
110985 & Antethic & Cheliuskin & Rock \\
129973 & Pipe Choir & Savior & Rock \\
136707 & Tojamura 7 & 100ft Wave & Rock \\
138060 & Every Now and Every Then & One Knife & Rock \\
146478 & Norwegian Blue Records & SPARKLE BLOOD - Denim Convention & Rock \\
148537 & E A Terra Nunca Me Pareceu Tão\ldots & Medo de Morrer & Rock \\
\midrule
\end{longtable}
\end{footnotesize}



\end{document}
